\documentclass[12pt, a4paper]{article}
\usepackage[utf8]{inputenc}
\usepackage[T1]{fontenc}
\usepackage{amsmath, amssymb, amsthm}
\usepackage{graphicx}
\usepackage{booktabs}
\usepackage{array}
\usepackage{multirow}
\usepackage{hyperref}
\usepackage{cite}
\usepackage{footnote}
\usepackage{geometry}
\usepackage{caption}
\usepackage{subcaption}
\usepackage{algorithm}
\usepackage{algpseudocode}
\usepackage{xcolor}
\usepackage{lipsum}
\usepackage{longtable}
\geometry{margin=2.5cm}

\providecommand{\keywords}[1]{\vspace{0.5em}\noindent\textbf{Keywords:} #1}

\title{Rainfall Estimation and Forecasting from Commercial Microwave Link Networks:\\
A Machine Learning Approach Using Swedish Cellular Infrastructure}

\author{\textit{Ben Yehoshua, S.}, \textit{Jacoby, D.}, \textit{Salganik, Y.}}

\date{\textit{Tel Aviv University} \\ \today}


\begin{document}
\maketitle

\begin{abstract}
Accurate quantitative precipitation estimation and short-term forecasting
are foundational requirements for urban hydrology, flash-flood early
warning, and operational meteorology. Conventional infrastructure relies
on rain gauges, which provide accurate but spatially sparse point
measurements, and weather radar, which provides spatially dense fields but
suffers from beam blockage, Z--R uncertainty, and high capital cost.
Commercial Microwave Links (CMLs), the backhaul of cellular networks,
constitute a dense and globally deployed opportunistic sensing
infrastructure: the rain-induced attenuation of these radio signals
encodes path-averaged rain rate. We present a complete pipeline that
converts raw CML received-signal-level (RSL) measurements into spatial
rainfall fields and short-term forecasts. Using the public openMRG
dataset (Gothenburg, Sweden, summer 2015), we build inverse-distance
weighting (IDW) and Goldshtein--Messer--Zinevich (GMZ) tomographic maps
from \(\sim\!214\) CMLs and benchmark a CML neural estimation network
against the ITU-R physical baseline. We further train three modern
sequence models---Transformer, Sparse Identification of Nonlinear
Dynamics (SINDy), and a Mamba-style state-space surrogate---in four
variants (IDW/GMZ \(\times\) single/multi-step) for nowcasts at horizons
of 15, 30, 45, and 60 minutes. Naive persistence baselines and the
operational pySTEPS Lagrangian extrapolator are used as references.
Performance will be reported using RMSE, MAE, bias, Pearson correlation,
and the Heidke Skill Score, evaluated against both SMHI rain gauges
(gold standard) and SMHI radar maps. Quantitative results will be
reported in a follow-up study upon completion of all model training.
\end{abstract}

\keywords{Commercial microwave links; rainfall estimation; precipitation nowcasting;
machine learning; Transformer; SINDy; Mamba; pySTEPS; opportunistic sensing;
Sweden; openMRG}

\section{Introduction}\label{sec:intro}

Quantitative precipitation estimation (QPE) and short-term forecasting
(``nowcasting'', the 0--2~hour horizon) underpin a broad range of
applications: urban drainage and flash-flood early warning, agriculture,
aviation safety, and the assimilation of observations into numerical
weather prediction (NWP) models. The two operational measurement
infrastructures---rain-gauge networks and weather radar---each present
well-known limitations. Rain gauges provide accurate, in~situ
measurements of accumulated rainfall but only at a single point, so that
dense urban or sparse rural networks alike fail to resolve the
considerable sub-kilometre variability of convective rainfall. Weather
radar provides high-resolution spatial fields but is affected by beam
blockage at low elevations, range-dependent overshooting of low cloud,
ground clutter, and the substantial uncertainty of the empirical
reflectivity-to-rain-rate relation (Z--R) used to convert radar
reflectivity into rain rate; deployment and maintenance costs are also
high.

The opportunistic sensing paradigm offers a complementary information
source. Commercial Microwave Links (CMLs)\footnote{CMLs are the
point-to-point microwave radio links that interconnect cellular base
stations. They typically operate in frequency bands between 15 and
40~GHz, where the wavelength is comparable to the size of raindrops; in
this regime, rainfall along the propagation path causes a measurable
attenuation of the received signal level (RSL) that is approximately
proportional to the path-integrated rain rate.} form a dense,
already-deployed, and continuously operating wireless infrastructure
whose attenuation measurements correlate strongly with the
path-integrated rain rate along each link. The same signal degradation
that operators monitor for quality-of-service purposes thus becomes a
rainfall sensor of opportunity at essentially zero marginal hardware cost.

The seminal demonstration by \mbox{Messer et al.} \cite{messer2006} that
country-scale rainfall fields can be inferred from cellular backhaul
attenuations triggered an active line of research. Closed-form physical
estimators based on the ITU-R~P.838 power law \cite{ituR838} were refined
by \mbox{Goldshtein et al.} \cite{goldshtein2009} into practical rain-rate
inversions, and the path-averaged nature of CML observations naturally
suggested tomographic reconstruction. National-scale operational
demonstrations \cite{overeem2013} followed. More recently, the
availability of public reference datasets and the maturity of deep
learning have motivated learning-based per-link rainfall estimators and
spatial reconstructions.

This study presents an end-to-end machine-learning pipeline anchored on
the public openMRG dataset \cite{andersson2023openmrg}, which combines
cellular microwave link data, SMHI rain gauges, and SMHI radar for
Gothenburg, Sweden, summer 2015. The contributions of this work are: (i)
an end-to-end pipeline that converts raw CML RSL into 15-minute
spatial rainfall maps using both IDW and GMZ reconstructions; (ii) a
trained CML estimation network that supplies the input maps to the
forecasting stage; and (iii) a systematic comparison of three modern
sequence models---Transformer, SINDy, and a Mamba-style state-space
surrogate---against naive persistence baselines and the operational
pySTEPS nowcast, at horizons of 15, 30, 45 and 60~min.

The remainder of the paper is organized as follows. Section
\ref{sec:background} reviews CMLs as rain sensors, the relevant
attenuation physics, spatial interpolation, nowcasting, and the radar
Z--R relation. Section~\ref{sec:dataset} describes the openMRG dataset
and our preprocessing. Section~\ref{sec:methodology} details every model.
Section~\ref{sec:metrics} defines the evaluation framework, and
Section~\ref{sec:setup} reports the experimental setup. Sections
\ref{sec:results} through \ref{sec:conclusion} present the planned
results, a forward-looking discussion, and concluding remarks.

\section{Background and Related Work}\label{sec:background}

\subsection{Commercial Microwave Links as Rain Sensors}

The use of cellular backhaul attenuation for environmental monitoring was
first proposed and demonstrated by \mbox{Messer et al.} \cite{messer2006}
in a landmark \emph{Science} paper that showed how the existing
infrastructure of a cellular operator could provide rainfall information
at country scale. Subsequent quantitative work by \mbox{Goldshtein et
al.} \cite{goldshtein2009} formalised the per-link inversion problem and
provided practical baseline-tracking algorithms, while \mbox{Overeem et
al.} \cite{overeem2013} produced country-wide rainfall maps for the
Netherlands using thousands of links from a single operator. The
underlying mechanism is the interaction of electromagnetic waves with
hydrometeors: at the frequencies used by cellular backhaul (typically
15--40~GHz), Mie scattering and absorption by raindrops produce a
specific attenuation that is well described by the
power law of ITU-R~Recommendation~P.838 \cite{ituR838}.

\subsection{Rain Attenuation Physics}\label{sec:physics}

The specific attenuation \(\gamma\) [dB/km] induced by rainfall on a
microwave link is related to the surface rain rate \(R\) [mm/h] by the
ITU-R power law:
\begin{equation}
  \gamma \;=\; a \cdot R^{\,b}, \label{eq:itu_power}
\end{equation}
where the coefficients \(a\) and \(b\) depend on the carrier
frequency, polarisation, and (weakly) drop-temperature; their values
are tabulated in \cite{ituR838} and reproduced in
Table~\ref{tab:itu_coeff}. For a link of physical length \(L\) [km], the
total path attenuation \(A\) [dB] satisfies
\begin{equation}
  A \;=\; \gamma \cdot L \;=\; a \cdot R^{\,b} \cdot L. \label{eq:itu_total}
\end{equation}
Inverting Equation~\eqref{eq:itu_total} yields the closed-form rain
estimator
\begin{equation}
  R \;=\; \left(\frac{A}{a\,L}\right)^{1/b}, \label{eq:itu_inverse}
\end{equation}
which constitutes the simplest possible CML rainfall estimator and the
physical baseline used in this study.

In practice, only the received signal level (RSL) of the link is
recorded; the total attenuation must be inferred from
\begin{equation}
  A(t) \;=\; \mathrm{RSL}_{\text{dry}}(t) \;-\; \mathrm{RSL}(t), \label{eq:rsl_to_atten}
\end{equation}
where \(\mathrm{RSL}_{\text{dry}}(t)\) denotes the level the link
\emph{would} have shown under dry conditions. Because antennas drift
thermally and link conditions change slowly over time, the dry baseline
is itself unknown and must be estimated from the data. We follow the
standard rolling-quantile approach
\begin{equation}
  \mathrm{RSL}_{\text{dry}}(t) \;=\; Q_{p}\!\left(
     \{\,\mathrm{RSL}(\tau) \;:\; \tau \in [t-W,\,t]\,\}\right), \label{eq:baseline}
\end{equation}
where \(Q_{p}\) is the \(p\)-th sample quantile and \(W\) is the window
length, exploiting the fact that for the great majority of the time it
is not raining and therefore the upper quantile of the RSL distribution
is a good estimate of the dry level. A common refinement subtracts a
small constant to correct for the so-called wet-antenna effect (WAE),
which is the residual attenuation caused by a film of water on the
antenna radome that persists for several minutes after a rain event
ends:
\begin{equation}
  A_{\text{corrected}}(t) \;=\; A(t) \;-\; A_{\mathrm{WAE}}. \label{eq:wae}
\end{equation}

\begin{table}[h]
\centering
\caption{Representative ITU-R~P.838 coefficients $(a, b)$ for selected
carrier frequencies and both horizontal (H) and vertical (V) linear
polarisations. The values shown are illustrative reference values taken
from the ITU-R~P.838-3 standard \cite{ituR838} and \emph{not} parameters
chosen in our codebase; in practice, \texttt{pynncml}'s
\texttt{one\_step\_dynamic\_baseline} routine looks up the full
ITU-R~P.838 table internally for each link based on its actual carrier
frequency and polarisation as recorded in the openMRG metadata.}
\label{tab:itu_coeff}
\begin{tabular}{lcccc}
\toprule
Frequency [GHz] & $a_{\mathrm H}$ & $b_{\mathrm H}$ & $a_{\mathrm V}$ & $b_{\mathrm V}$ \\
\midrule
15 & 0.0376 & 1.154 & 0.0335 & 1.128 \\
18 & 0.0691 & 1.099 & 0.0691 & 1.099 \\
23 & 0.1284 & 1.051 & 0.1284 & 1.051 \\
26 & 0.1402 & 1.061 & 0.1402 & 1.061 \\
38 & 0.4001 & 0.929 & 0.3950 & 0.913 \\
\bottomrule
\end{tabular}
\end{table}

\subsection{Spatial Interpolation Methods}\label{sec:interp}

Both per-link rain rates and gauge readings are scattered (line or
point) measurements; reconstructing a regular field requires spatial
interpolation. We use two complementary methods.

\paragraph{Inverse Distance Weighting (IDW).} Each grid point
\(\mathbf{x}\) is estimated as a weighted mean of the source values
\(z_i\) at locations \(\mathbf{x}_i\):
\begin{equation}
  \hat{z}(\mathbf{x}) \;=\;
  \frac{\sum_{i=1}^{N} w_i(\mathbf{x})\, z_i}
       {\sum_{i=1}^{N} w_i(\mathbf{x})},
  \qquad
  w_i(\mathbf{x}) \;=\; \frac{1}{d(\mathbf{x},\mathbf{x}_i)^{p}},
  \label{eq:idw}
\end{equation}
with \(p=2\) and Euclidean distance evaluated in lon/lat (degrees);
this approximation is sufficient over the small Gothenburg domain. IDW
is computationally cheap and stable but it is a known limitation of
the method that intensity peaks are systematically smoothed toward the
local mean: because every grid point is a convex combination of nearby
measurements, the interpolated surface cannot exceed the local maximum
or fall below the local minimum of the observation set, and rapidly
varying features---such as the cores of convective cells---are
attenuated relative to the true field. The effect has been quantified
specifically in the CML rainfall-mapping context by Eshel et al.\
\cite{eshel2021spatial}, who showed that IDW underestimates rainfall
peaks compared to methods that exploit the path-integral nature of CML
observations.

\paragraph{Goldshtein--Messer--Zinevich (GMZ) tomography.}
The GMZ reconstruction treats each link's path attenuation as a path
integral of the (unknown) two-dimensional rain field and solves for the
field that best reproduces all observed CML attenuations
simultaneously. The forward model along link \(\ell\) of length
\(L_\ell\) discretises the path integral with \(N\) sample points:
\begin{equation}
  \bar{R}_\ell \;=\;
  \frac{1}{L_\ell}\int_0^{L_\ell}\! R(s)\, ds
  \;\approx\;
  \frac{1}{N}\sum_{j=1}^{N} R(s_j). \label{eq:gmz_path}
\end{equation}
The original algorithm is due to Goldshtein, Messer and Zinevich
\cite{goldshtein2009}; subsequent CML mapping studies typically
represent each link by a small number of virtual rain gauges (VRGs)
distributed along its path \cite{eshel2021spatial}, recovering a
sharper estimate of the rain field than IDW at the cost of additional
computation. The \texttt{pynncml} toolbox
(\texttt{mcm.generate\_link\_set\_gmz}) exposes the two engineering
parameters that control this discretisation:
\texttt{region\_of\_interest}, the pixel-radius of each link's region
of influence on the reconstruction grid, and \texttt{point\_per\_link},
the number of VRG sample points used to discretise the path integral
of Equation~\eqref{eq:gmz_path}. These parameters are
implementation-specific to the \texttt{pynncml} GMZ routine and are
not, to our knowledge, the subject of dedicated published guidance;
the literature uses comparable but differently-named discretisations
(e.g.\ Eshel et al.\ \cite{eshel2021spatial} compare 1-VRG and 3-VRG
representations). For consistency with internal prior work in our
group, we fix \(\texttt{REGION\_OF\_INTEREST}=3\) pixels and
\(\texttt{POINT\_PER\_LINK}=2\) VRGs across all experiments\footnote{We
note that these two values are engineering choices: they were neither
optimised against the openMRG validation split nor inherited from a
published reference. They were carried over from earlier internal
experiments so that all results in this study remain directly
comparable to prior work conducted in the laboratory, and a dedicated
study of their sensitivity is left to future work.}.

\subsection{Precipitation Nowcasting}

Nowcasting designates short-term (0--2~hour) forecasting of precipitation.
At these horizons, the bulk of the forecast skill derives from the
extrapolation of observed precipitation features rather than from the
explicit modelling of atmospheric dynamics. The classical approach
estimates a motion vector field from a sequence of recent radar maps
and advects the most recent observation forward (Lagrangian
extrapolation):
\begin{equation}
  R(\mathbf{x},\, t+h) \;\approx\; R(\mathbf{x} - \mathbf{v}\cdot h,\; t), \label{eq:advection}
\end{equation}
where \(\mathbf{v}\) is the (possibly spatially varying) motion vector
field. The open-source pySTEPS library \cite{pulkkinen2019pysteps}
implements a state-of-the-art version of this scheme with
Lucas--Kanade or variational optical-flow motion estimators and an
optional stochastic ensemble for probabilistic forecasting. Deep
learning has more recently been used to learn the advection plus growth
and decay jointly from data, but the deterministic pySTEPS pipeline
remains the operational reference against which new methods are
typically benchmarked.

\subsection{Radar Reflectivity and Z--R Relations}

Weather radars do not measure rain rate directly: they measure
backscattered power, summarized as the radar reflectivity factor
\(Z\) [mm\(^6\)/m\(^3\)]. The connection to rain rate is the empirical
power law
\begin{equation}
  Z \;=\; a_Z \cdot R^{\,b_Z}, \label{eq:zr}
\end{equation}
whose canonical Marshall--Palmer parameters \cite{marshallpalmer1948}
are \(a_Z=200,\; b_Z=1.6\). Inverting Equation~\eqref{eq:zr} with the
common dBZ scaling \(Z_{\text{dBZ}}=10\log_{10}Z\) gives
\begin{equation}
  R \;=\; \left(\frac{10^{\,Z_{\text{dBZ}}/10}}{a_Z}\right)^{1/b_Z}. \label{eq:zr_inv}
\end{equation}
This relation is used in our preprocessing pipeline to convert the
openMRG radar dBZ field into rain rate.

\section{Dataset}\label{sec:dataset}

\subsection{The openMRG Dataset}

We use the public \emph{Open Microwave Rainfall Gothenburg} (openMRG)
dataset \cite{andersson2023openmrg}, which co-locates three precipitation
data streams over the Gothenburg area, Sweden, during the summer of
2015: Commercial Microwave Links, rain gauges, and SMHI weather radar.
The native openMRG CML inventory comprises 364 bi-directional links
(728 sub-links in total), with path lengths from approximately 100~m
to 15~km (median \(\approx 2\) km), carrier frequencies in the range
7--38~GHz (with the majority above 25~GHz), and predominantly vertical
polarisation. The CMLs are instrumented with min/max received and
transmitted signal levels (RSL and TSL) aggregated at a native temporal
resolution of 10~seconds, with reported coverage of 99.99\% over the
study period \cite{andersson2023openmrg}. The full openMRG period is
1~June 2015 to 31~August 2015. Per-link metadata include the two site
coordinates, the path length, the carrier frequency, and the
polarisation. The dataset is loaded via the \texttt{pynncml} toolbox,
which provides a unified loader (\texttt{loader\_open\_mrg\_dataset})
and dataset object that bundles the link set, the gauge point set, and
the time alignment used throughout this study; after \texttt{pynncml}'s
internal quality filtering and the selection of a single direction per
bi-directional pair, approximately 214 link series are exposed for
use in this work.

\begin{table}[h]
\centering
\caption{Summary of the openMRG dataset components used in this study.}
\label{tab:openmrg_summary}
\begin{tabular}{lll}
\toprule
Component & Quantity & Notes \\
\midrule
CMLs (raw)                  & 364 bi-directional             & 728 sub-links; 7--38 GHz, mostly vertical pol. \\
CMLs (this work)            & \(\sim\)214 series             & after \texttt{pynncml} filtering / direction pick \\
CML temporal resolution     & 10 s                            & min/max RSL/TSL \\
CML path length             & 0.1--15 km                      & median \(\approx 2\) km \\
Rain gauges                 & 11 stations                    & 1 SMHI weighing + 10 city network \\
Gauge temporal resolution   & 15 min (SMHI) / 1 min (city)   & \(\sim\)10 city gauges used here \\
Radar (SMHI composite)      & 1 field                        & 48\(\times\)37 pixels, 5-min cadence \\
Radar reflectivity encoding & pseudo-dBZ (uint8)             & \(\mathrm{dBZ} = 0.4\!\cdot\!\text{raw} - 30\) \\
Time span                   & 2015-06-01 to 2015-08-31       & 92 days \\
Crop domain (lat)           & 57.55--57.85 \(^\circ\)N       & \\
Crop domain (lon)           & 11.70--12.20 \(^\circ\)E       & \\
Working temporal resolution & 15 minutes                     & via right-aligned-mean resampling \\
\bottomrule
\end{tabular}
\end{table}

\subsection{SMHI Rain Gauge Network}\label{sec:gauges}

Rain gauges are point measurements of precipitation at a single
location, traditionally regarded as the gold standard for surface
rainfall over the area immediately surrounding the
station\footnote{Although gauges are spatially limited to a single
point and may underestimate intense rainfall due to wind-induced
under-catch and instrumental losses, the directness of the measurement
(physical collection of water in a known catchment area, followed by a
weighing or tipping mechanism) yields an accuracy that no remote
sensing technique currently matches. Gauges are therefore the customary
reference against which CML and radar estimates are validated.}.
The openMRG bundle provides eleven gauges in total: a single SMHI
automatic weather station equipped with a Geonor weighing gauge
(0.1~mm resolution, native 15-min cadence) plus a ten-station city
network operated by Gryaab (seven Geonor weighing and three
tipping-bucket gauges, the latter at 0.2~mm resolution, all at a
1-minute cadence). The pynncml dataloader exposes the high-cadence
city-network stations as its \texttt{point\_set}, giving the
\(\sim\)10 gauges used as the ground-truth reference throughout this
study. The 1-minute rain-rate series are resampled to the radar's
15-min cadence using a right-aligned mean and aligned with the radar
time grid via nearest-neighbour matching within a tolerance of
8 minutes.

\subsection{SMHI Weather Radar}

The openMRG bundle includes a single SMHI two-dimensional radar
composite field over the Gothenburg domain, at approximately 2~km
horizontal resolution and 5-minute temporal resolution; the openMRG
file contains a 48\(\times\)37-pixel field, of which the cropped
sub-window of Table~\ref{tab:openmrg_summary} is retained. The
composite is built from the national SMHI radar network (with the
nearest contributing radars being Vara, Karlskrona, and Vilebo, of
which Vara is closest at \(\approx 78\) km), so the Gothenburg domain
is observed at significant range\footnote{The radar limitations
discussed below are exacerbated by this range: the nominal vertical
sampling of the radar beam is several hundred meters above the surface
at such distances, so beam overshooting of low precipitation and
range-dependent under-catch are intrinsic limitations of the openMRG
radar product over the study area.}. The raw field is encoded as an
unsigned 8-bit pseudo-dBZ representation; values of 255 indicate
missing data and the remaining values are scaled via
\(\mathrm{dBZ} = 0.4\!\cdot\!\text{raw} - 30\)
\cite{andersson2023openmrg}. The radar serves as the secondary spatial
reference in this study\footnote{Weather radar is
spatially dense and operationally important, but the underlying Z--R
relationship is uncertain and varies with precipitation regime, drop
size distribution and viewing geometry; range degradation, beam
blockage and bright-band contamination further bias the rainfall
estimate. Quantitative validation of CML methods is therefore
preferentially done against gauges, with radar used as a spatial
reference rather than a ground truth.}.

\subsection{Data Preprocessing}\label{sec:preproc}

The raw streams are aligned onto a common 15-minute grid and processed
as follows.

\paragraph{Radar.} Raw uint8 values of 255 are masked as missing, the
remainder is converted to dBZ via \(\mathrm{dBZ} = \text{raw}\cdot 0.4
- 30\). To suppress hail-induced reflectivities the field is capped at
53~dBZ, then converted to rain rate via the Marshall--Palmer inversion
of Equation~\eqref{eq:zr_inv} with \(a_Z=200,\; b_Z=1.6\). The result
is resampled to 15-min cadence with a right-aligned mean and stored as
a \((T, H, W)\) array.

\paragraph{CML.} The native openMRG CML stream provides min/max RSL
and TSL aggregated over 10-second intervals
\cite{andersson2023openmrg}; \texttt{pynncml}'s OpenMRG dataloader
ingests these series and exposes per-link minute-cadence tensors that
feed the downstream estimators. The rolling-quantile dry baseline of
Equation~\eqref{eq:baseline} is applied implicitly inside
\texttt{pynncml}'s \texttt{one\_step\_dynamic\_baseline} routine. A
wet/dry decision uses a wet threshold of 0.5~dB; baseline window is
60~min, and a minimum reportable rain rate of \(R_{\min} = 0.1\) mm/h
is enforced. The per-link rain rates are finally resampled to 15-min
cadence and aligned with the radar grid via nearest-neighbour matching
with a tolerance of 8 minutes.

\paragraph{Quality filtering.} Links with non-physical or excessively
NaN-laden series are excluded by \texttt{pynncml}'s dataloader; the
remaining \(\sim\)214 links cover the full openMRG period.

\paragraph{Train/validation/test split.} A strictly temporal split is
used throughout: all training, validation, and testing series obey the
chronological boundaries in Table~\ref{tab:splits}. No random
shuffling of timestamps is performed at any stage.

\begin{table}[h]
\centering
\caption{Chronological data split used for all model training and evaluation.}
\label{tab:splits}
\begin{tabular}{lll}
\toprule
Split & Date range & Purpose \\
\midrule
Train & 2015-06-01 to 2015-07-20 & Model fitting (50 days) \\
Val   & 2015-07-21 to 2015-08-05 & Hyperparameter selection, early stopping (16 days) \\
Test  & 2015-08-06 to 2015-08-31 & Held-out evaluation (26 days) \\
\bottomrule
\end{tabular}
\end{table}

\section{Methodology}\label{sec:methodology}

\subsection{Pipeline Overview}

The end-to-end pipeline (Fig.~\ref{fig:pipeline}) proceeds in four
stages. Raw per-link RSL/TSL series are converted into per-link
attenuations via dynamic baseline estimation
(Eq.~\ref{eq:rsl_to_atten}--\ref{eq:wae}). A rain-rate estimator
(either the physical ITU-R inversion or a learned neural network)
produces a per-link rain rate time series. The scattered per-link
estimates are projected onto the regular radar grid via IDW
(Eq.~\ref{eq:idw}) and GMZ (Eq.~\ref{eq:gmz_path}), producing
two-dimensional rainfall map sequences at 15-min cadence. Finally,
forecast models---naive persistence, Transformer, SINDy, Mamba and the
operational pySTEPS---ingest sequences of these maps and produce
forecast maps at horizons of 15, 30, 45 and 60~min.

\begin{figure}[h]
\centering
\fbox{\textit{[Pipeline diagram --- to be inserted]}}
\caption{End-to-end pipeline from raw CML measurements to rainfall
forecast maps. Stage 1: per-link attenuation extraction with dynamic
baseline. Stage 2: per-link rain rate (physical or learned). Stage 3:
spatial reconstruction (IDW or GMZ). Stage 4: short-term forecasting.}
\label{fig:pipeline}
\end{figure}

\subsection{Model 1 --- Physical Baseline (ITU-R Inversion)}\label{sec:model1}

Model~1 is the direct, learning-free inversion of the ITU-R power law
(Eq.~\ref{eq:itu_inverse}). The implementation reuses
\texttt{pynncml.scm.rain\_estimation.one\_step\_dynamic\_baseline} with
the \texttt{MAX} power-law variant, a minimum reportable rain rate of
\(R_{\min} = 0.1\) mm/h, a dynamic baseline window of 60~min, and a
wet/dry detection threshold of 0.5~dB. The per-link rain-rate series
is then projected to the radar grid via IDW
(\(p=2\)) to produce \texttt{model1\_physical}, the physical baseline
field. No learning is involved; the only choices are the baseline
window, the wet/dry threshold, and the interpolation method.

\subsection{Model 2 --- CML Neural Estimation Model}\label{sec:model2}

Model~2 is the learned counterpart to Model~1. We use the two-step
network architecture from \texttt{pynncml.scm.rain\_estimation.two\_step\_network},
which consists of a multi-layer GRU recurrent encoder followed by two
parallel output heads: a rain-rate regressor and a wet/dry binary
detector. The detector probability gates the regressor at inference
time, which suppresses spurious rain rates on dry samples. The input
to the network at each time step is the concatenation of the
\emph{instantaneous} RSL and TSL values for a single link together
with two metadata features (link length and carrier frequency); the
target is the corresponding gauge-derived per-link rain rate. The
network is trained on temporal windows of length
\(\texttt{WINDOW\_SIZE} = 32\) one-minute samples.

The composite training loss combines a weighted mean-squared error for
the rain-rate head with a binary cross-entropy term for the detection
head. The MSE weighting follows the
\texttt{RegressionLoss} formulation
\begin{equation}
  \mathcal{L}_{\text{reg}}(\hat{r}, r) =
  \sum_{t}\frac{1}{B}\sum_{b}\Bigl(1 - \gamma_s\,e^{-\gamma\,r_{b,t}}\Bigr)\,
                (r_{b,t} - \hat{r}_{b,t})^2, \label{eq:reg_loss}
\end{equation}
which up-weights samples with non-zero rain, since the data are
strongly imbalanced toward dry. The detection head uses standard
binary cross-entropy against the indicator \(\mathbf{1}\{r > 0.1\}\).
The total loss is
\(\lambda\,\mathcal{L}_{\text{reg}} + \mathcal{L}_{\text{det}}\), with
\(\lambda\) calibrated automatically from the first training batch.

Hyperparameters are selected by a Bayesian W\&B sweep
(project \texttt{FieldSense-Estimation-v2}, 10 trials, 30 epochs each)
over the learning rate (log-uniform \(10^{-5}\)--\(10^{-3}\)), batch
size (\{8, 16, 32\}), RNN hidden dimensionality (\(\{128, 256\}\)),
number of GRU layers (\(\{1, 2\}\)), and weight decay (log-uniform
\(10^{-6}\)--\(10^{-3}\)); the optimiser is RAdam with gradient
clipping at \(\lVert g\rVert_2 \le 1.0\). The best-by-validation-loss
configuration is retrained for \(\texttt{FINAL\_EPOCHS\_EST} = 200\)
epochs with best-by-validation checkpointing. After training, the
per-link rain rates predicted by Model~2 are projected to the radar
grid via both IDW and GMZ to produce the two-dimensional
\texttt{model2\_cml\_estimation\_maps\_idw} and
\texttt{\_gmz} arrays, which become the input maps for all forecast
models in Sections~\ref{sec:tf}--\ref{sec:mamba}.

\begin{table}[h]
\centering
\caption{Architecture of the CML estimation network (Model~2). The
recurrent dimensionality and layer count are W\&B-tuned within the
indicated search space; the values shown in parentheses are the lab's
reference defaults.}
\label{tab:model2_arch}
\begin{tabular}{llll}
\toprule
Layer & Type & Output dim. & Activation \\
\midrule
Input  & RSL, TSL, metadata           & \(2 + 2\)                  & --- \\
Encoder & GRU, \texttt{n\_layers} (\(\{1,2\}\), ref.~2) & \texttt{rnn\_n\_features} (\(\{128,256\}\), ref.~256) & tanh \\
Head 1 (reg.) & Linear $\rightarrow$ Linear & 1 & Softplus \\
Head 2 (det.) & Linear $\rightarrow$ Linear & 1 & Sigmoid \\
\bottomrule
\end{tabular}
\end{table}

\subsection{Naive Persistence Models}\label{sec:naive}

Persistence---``the world does not change''---is the trivial forecast
\begin{equation}
  \hat{Y}(t+h) \;=\; Y(t), \label{eq:persist}
\end{equation}
and it is a famously hard baseline to beat at short horizons because
rain fields are strongly autocorrelated in time. We build three
variants of Equation~\eqref{eq:persist}, distinguished by the source
field \(Y\):
\begin{align}
  \hat{M}_{\text{radar}}(t+h) &= M_{\text{radar}}(t), \label{eq:naive_radar} \\
  \hat{M}_{\text{gauge}}(t+h) &= \mathrm{IDW}\!\left(\{r_s(t)\}_{s=1}^{S}\right), \label{eq:naive_gauge} \\
  \hat{M}_{\text{CML}}(t+h)   &= \hat{M}_{2}(t), \label{eq:naive_cml}
\end{align}
with \(h \in \{15, 30, 45, 60\}\) minutes. Equation~\eqref{eq:naive_radar}
applies persistence to the radar field directly,
Equation~\eqref{eq:naive_gauge} to the IDW-interpolated gauge field, and
Equation~\eqref{eq:naive_cml} to the Model-2 IDW estimation map. These
baselines provide an interpretable floor for every learned forecaster.

\subsection{Common Forecasting Dataset Construction}\label{sec:fcst_data}

All learned forecast models share the same supervised data layout.
Given a \((T, H, W)\) sequence of 15-min maps from one of the
\texttt{MAP\_SOURCES} (IDW or GMZ) and a fixed lookback
\(\texttt{LOOKBACK\_WINDOW} = w = 8\) frames (\(=2\) hours of
history), we form supervised pairs
\begin{align}
  X_i &= \bigl(M_{i},\, M_{i+1},\, \ldots,\, M_{i+w-1}\bigr) \;\in\; \mathbb{R}^{w\times H\times W}, \\
  Y_i &= M_{i+w+h-1} \;\in\; \mathbb{R}^{H\times W} \qquad \text{(single-step)},
\end{align}
or, for multi-step training,
\begin{equation}
  Y_i \;=\; \bigl(M_{i+w},\, M_{i+w+1},\, M_{i+w+2},\, M_{i+w+3}\bigr) \;\in\; \mathbb{R}^{n_h\times H\times W},
\end{equation}
where \(h \in \{1, 2, 3, 4\}\) corresponds to lead times of
\(15h\) minutes and \(n_h = 4\). NaNs are zeroed before
tensorisation. Splits use the dates of Table~\ref{tab:splits}.

\subsection{Transformer-Based Forecasting}\label{sec:tf}

The Transformer \cite{vaswani2017attention} relies on the scaled
dot-product self-attention operation
\begin{equation}
  \mathrm{Attention}(\mathbf{Q}, \mathbf{K}, \mathbf{V}) \;=\;
  \mathrm{softmax}\!\left(\frac{\mathbf{Q}\mathbf{K}^{\top}}{\sqrt{d_k}}\right)\mathbf{V},
  \label{eq:attention}
\end{equation}
where each row of \(\mathbf{Q}\in\mathbb{R}^{T\times d_k}\) queries a
collection of key--value pairs and the softmax produces a
context-dependent weighting over time. To preserve order information,
sinusoidal positional encodings are added to the token embeddings:
\begin{equation}
  \mathrm{PE}(pos, 2i) = \sin\!\left(\frac{pos}{10000^{\,2i/d_{\text{model}}}}\right), \quad
  \mathrm{PE}(pos, 2i+1) = \cos\!\left(\frac{pos}{10000^{\,2i/d_{\text{model}}}}\right). \label{eq:pe}
\end{equation}
Multi-head attention concatenates \(h_a\) independent attention
operations of dimension \(d_k = d_{\text{model}}/h_a\) and applies a
final linear projection.

Our \texttt{RainTransformer} flattens each \((H, W)\) input frame into
a single token of dimension \(H\cdot W = 1{,}776\), projects it to the
model dimension \(d_{\text{model}}\) by a linear embedding layer,
processes the resulting sequence of \(w\) tokens with a stack of
\texttt{n\_layers} \texttt{nn.TransformerEncoderLayer} blocks
(multi-head attention with \texttt{n\_heads} heads followed by a
feed-forward sublayer), and then decodes the last-token
representation through a single linear layer that outputs the next
\(n_h\) frames flattened. Single-step models output one frame; the
multi-step model outputs all four horizons in a single forward pass.
A separate model is trained for each of the four 2\(\times\)2 variants
(IDW/GMZ \(\times\) single/multi-step), giving eight Transformer
checkpoints in total.

Hyperparameters are selected by a Bayesian W\&B sweep (project
\texttt{FieldSense-Transformer-v2}, \(\texttt{SWEEP\_COUNT\_TF} = 6\)
trials of \(\texttt{SWEEP\_EPOCHS\_TF} = 25\) epochs each, single-step
horizon-1 IDW target) over the parameters in
Table~\ref{tab:tf_sweep}; the best configuration is saved to
\texttt{transformer\_sweep\_best.json} and reused for all eight final
trainings (\(\texttt{FINAL\_EPOCHS\_TF} = 80\) epochs with
best-by-validation checkpointing, Adam optimiser, MSE loss, gradient
clipping at \(\lVert g\rVert_2 \le 1\)).

\begin{table}[h]
\centering
\caption{Transformer hyperparameter search space.}
\label{tab:tf_sweep}
\begin{tabular}{lll}
\toprule
Hyperparameter & Search space & Type \\
\midrule
\(d_{\text{model}}\) & \(\{64, 128, 256\}\)              & categorical \\
\(n_{\text{heads}}\) & \(\{2, 4, 8\}\)                   & categorical \\
\(n_{\text{layers}}\) & \(\{1, 2, 3\}\)                  & categorical \\
learning rate & \(\mathcal{U}_{\log}[10^{-4}, 10^{-3}]\) & log-uniform \\
batch size    & \(\{8, 16, 32\}\)                        & categorical \\
\bottomrule
\end{tabular}
\end{table}

\subsection{SINDy-Based Forecasting}\label{sec:sindy}

Sparse Identification of Nonlinear Dynamics
(SINDy)\footnote{Following Brunton, S.L., Proctor, J.L., \& Kutz, J.N.
(2016). \emph{Discovering governing equations from data by sparse
identification of nonlinear dynamical systems.} PNAS,
113(15):3932--3937.} \cite{brunton2016sindy} discovers a parsimonious
system of ordinary differential equations directly from data. Given a
matrix of state snapshots \(\mathbf{X}\in\mathbb{R}^{T\times n}\) and
their numerical time derivatives \(\dot{\mathbf{X}}\), SINDy fits
\begin{equation}
  \dot{\mathbf{X}} \;=\; \Theta(\mathbf{X})\,\boldsymbol{\Xi}, \label{eq:sindy_dyn}
\end{equation}
where the library matrix
\begin{equation}
  \Theta(\mathbf{X}) \;=\; \begin{bmatrix}
    \mathbf{1} & \mathbf{X} & \mathbf{X}^{2} & \cdots & \sin(\mathbf{X}) & \cdots
  \end{bmatrix} \label{eq:sindy_lib}
\end{equation}
collects candidate features (we use a polynomial library of degree~2)
and the coefficient matrix \(\boldsymbol{\Xi}\) is estimated by a
sparsity-promoting regression. The implementation in
\texttt{pysindy}'s \texttt{STLSQ} (sequentially thresholded least
squares with \(\ell_2\) ridge regularisation) solves
\begin{equation}
  \hat{\boldsymbol{\Xi}} \;=\;
  \arg\min_{\boldsymbol{\Xi}}
  \Bigl\|\dot{\mathbf{X}} - \Theta(\mathbf{X})\,\boldsymbol{\Xi}\Bigr\|_{2}^{2}
  \;+\; \lambda\,\|\boldsymbol{\Xi}\|_{1}, \label{eq:sindy_lasso}
\end{equation}
where we use \(\lambda = \texttt{SINDY\_THRESH} = 0.05\) and an
\(\ell_2\) coefficient \(\alpha = 0.05\). The discretisation timestep is
\(\Delta t = 0.25\) h (15~min). These three values, together with the
polynomial library degree, are held fixed at the lab's reference
configuration (see Table~\ref{tab:hparam_provenance}); only the number
of POD modes \(n_{\text{modes}}\) is data-driven via the W\&B sweep
described below.

Because the \(1{,}776\)-dimensional spatial state is too large for SINDy
to handle directly, we first reduce the state via truncated Singular
Value Decomposition (POD): \(\mathbf{X}_{\text{flat}} \in
\mathbb{R}^{T\times 1776}\) is projected onto \(n_{\text{modes}}\)
left-singular vectors via \texttt{TruncatedSVD}, and SINDy is fit on
the resulting POD coefficients. The number of modes,
\(n_{\text{modes}}\), is the most important SINDy hyperparameter and
is selected by a dedicated W\&B sweep (project
\texttt{FieldSense-SINDy-v2}) over \(n_{\text{modes}} \in
\{4, 6, 8, 10, 12, 16\}\); each setting is fitted on the training
maps, simulated forward at horizon~1 over the validation window, and
ranked by validation RMSE (and HSS for diagnostic purposes). The
best \(n_{\text{modes}}\) is then reused for all four SINDy variants
(IDW/GMZ \(\times\) single/multi-step). Forecasting at horizon \(h\)
is performed by integrating the discovered ODE system from the most
recent POD coefficient state forward to time \(h\Delta t\) using
\texttt{pysindy.SINDy.simulate} and inverse-projecting back to the
spatial grid.

\subsection{Mamba State Space Model}\label{sec:mamba}

State-space models (SSMs) describe a continuous-time linear dynamical
system of the form
\begin{equation}
  \mathbf{h}'(t) \;=\; \mathbf{A}\,\mathbf{h}(t) + \mathbf{B}\,x(t), \label{eq:ssm_state}
\end{equation}
\begin{equation}
  y(t) \;=\; \mathbf{C}\,\mathbf{h}(t), \label{eq:ssm_obs}
\end{equation}
where \(\mathbf{h}\) is the latent state, \(x\) the input, and
\(\mathbf{A}\), \(\mathbf{B}\), \(\mathbf{C}\) the (potentially
input-dependent) system matrices. The recent Mamba
architecture\footnote{Gu, A., \& Dao, T. (2023). \emph{Mamba:
Linear-time sequence modeling with selective state spaces.}
arXiv:2312.00752. The selective-scan mechanism allows the model to
focus or ignore information selectively based on the current input,
recovering input-dependence at the cost of forfeiting closed-form
convolutional implementations.} \cite{gu2023mamba} introduces a
\emph{selective} scan: \(\mathbf{A}\), \(\mathbf{B}\), and \(\mathbf{C}\)
become functions of the input \(x_t\), allowing the model to focus
attention on relevant past states. Discretisation to fixed step size
\(\Delta\) follows
\begin{equation}
  \bar{\mathbf{A}} \;=\; e^{\Delta\mathbf{A}},
  \qquad
  \bar{\mathbf{B}} \;=\; (\Delta\mathbf{A})^{-1}\bigl(e^{\Delta\mathbf{A}} - \mathbf{I}\bigr)\,\Delta\mathbf{B}.
  \label{eq:mamba_disc}
\end{equation}
The principal practical advantage of Mamba over the Transformer is
linear (rather than quadratic) complexity in sequence length.

We use a lightweight \texttt{SimpleMamba} surrogate that captures the
SSM recurrence pattern of Equation~\eqref{eq:ssm_state}--\eqref{eq:ssm_obs}
via a multi-layer GRU; this surrogate matches our group's earlier
internal reference implementation and avoids the additional engineering
overhead of a full selective-scan layer. The model flattens each
\((H, W)\) input frame to dimension \(H\cdot W = 1{,}776\), runs an
\texttt{nn.GRU} of size \(\texttt{hidden\_dim}\) and \(\texttt{n\_layers}\)
layers across the lookback, and decodes the final hidden state via a
single linear layer to \(n_h\cdot H\cdot W\) outputs. As for the
Transformer, eight checkpoints (IDW/GMZ \(\times\) single/multi-step) are
trained, with hyperparameters selected by a Bayesian W\&B sweep
(project \texttt{FieldSense-Mamba-v2}, \(\texttt{SWEEP\_COUNT\_MB} = 6\)
trials, \(\texttt{SWEEP\_EPOCHS\_MB} = 25\) epochs) over the search
space in Table~\ref{tab:mb_sweep}; the best configuration is reused for
\(\texttt{FINAL\_EPOCHS\_MB} = 80\) epochs of final training.

\begin{table}[h]
\centering
\caption{Mamba (SSM surrogate) hyperparameter search space.}
\label{tab:mb_sweep}
\begin{tabular}{lll}
\toprule
Hyperparameter & Search space & Type \\
\midrule
\(\texttt{hidden\_dim}\) & \(\{128, 256, 512\}\)             & categorical \\
\(\texttt{n\_layers}\)   & \(\{1, 2\}\)                      & categorical \\
learning rate            & \(\mathcal{U}_{\log}[10^{-4}, 10^{-3}]\) & log-uniform \\
batch size               & \(\{8, 16, 32\}\)                 & categorical \\
\bottomrule
\end{tabular}
\end{table}

\subsection{pySTEPS Nowcasting Baseline}\label{sec:pysteps}

pySTEPS \cite{pulkkinen2019pysteps} is used in its deterministic
extrapolation mode. For each test frame \(t\), the three preceding
radar frames are passed to the Lucas--Kanade optical-flow motion
estimator, producing a two-dimensional motion field \(\mathbf{v}\);
the current radar frame is then advected forward using
Equation~\eqref{eq:advection} for \(h \in \{1, 2, 3, 4\}\) steps of 15
minutes. Crucially, pySTEPS is fed the \emph{radar} field directly,
not the CML-derived maps, because it is designed for radar input; this
is therefore the operational state-of-the-art benchmark against which
the CML-based forecasters must measure themselves. The probabilistic
ensemble mode is not used in this study.

\subsection{Lookback Window Analysis}

The lookback \(w\) is a global hyperparameter shared by all learned
forecast models. A dedicated diagnostic study probes
\(\texttt{LOOKBACK\_VALUES} = \{2, 4, 6, 8, 10, 12, 16\}\) by training
a Transformer for \(\texttt{PROBE\_EPOCHS} = 10\) epochs on IDW maps
at horizon~1 and recording the validation MSE. The reference value of
\(w=8\) (2~hours) is retained throughout this paper; the probe serves
to verify that this choice is near-optimal and to inform future
reproducibility studies if the curve shows a clear minimum elsewhere.
Weights from the probe runs are discarded.

\subsection{Single-Step vs Multi-Step Prediction}\label{sec:singlemulti}

For each architecture we train two distinct prediction variants. The
single-step variant uses a separate model for each horizon,
\begin{equation}
  \hat{Y}_{h} \;=\; f_{h}\!\left(\mathbf{X}_{t-w:t}\right),
  \quad h \in \{15, 30, 45, 60\}, \label{eq:single}
\end{equation}
yielding \(n_h\) independent regressors. The multi-step variant uses a
single model that predicts all four horizons jointly:
\begin{equation}
  \bigl(\hat{Y}_{15}, \hat{Y}_{30}, \hat{Y}_{45}, \hat{Y}_{60}\bigr) \;=\; f\!\left(\mathbf{X}_{t-w:t}\right). \label{eq:multi}
\end{equation}
The single-step variant avoids error accumulation across horizons and
allows each model to specialise, at the cost of \(n_h\times\) the
parameter budget and the loss of shared representations. The
multi-step variant amortises the cost of feature extraction over all
horizons and benefits from multi-task regularisation, but couples
errors across the prediction heads. Both variants are evaluated.

\section{Evaluation Metrics}\label{sec:metrics}

All metrics are computed pixel-wise over the cropped Gothenburg domain
(or point-wise at gauge locations) and aggregated across the test
window. Let \(N\) denote the number of finite paired predictions and
observations and let \(\hat{y}_i\) and \(y_i\) denote the predicted and
observed values, respectively.

\paragraph{Magnitude error.} Root-mean-squared and mean-absolute
errors quantify the magnitude of pointwise residuals:
\begin{equation}
  \mathrm{RMSE} \;=\;
  \sqrt{\frac{1}{N}\sum_{i=1}^{N}\bigl(\hat{y}_i - y_i\bigr)^{2}}, \label{eq:rmse}
\end{equation}
\begin{equation}
  \mathrm{MAE} \;=\;
  \frac{1}{N}\sum_{i=1}^{N}\bigl|\hat{y}_i - y_i\bigr|. \label{eq:mae}
\end{equation}

\paragraph{Bias.} Systematic over- or under-estimation is summarised
by
\begin{equation}
  \mathrm{Bias} \;=\;
  \frac{1}{N}\sum_{i=1}^{N}\bigl(\hat{y}_i - y_i\bigr). \label{eq:bias}
\end{equation}

\paragraph{Linear correlation.} The Pearson product-moment correlation
\begin{equation}
  \mathrm{CC} \;=\;
  \frac{\mathrm{Cov}\!\left(\hat{y}, y\right)}
       {\sigma_{\hat{y}}\,\sigma_{y}}, \label{eq:cc}
\end{equation}
takes values in \([-1, 1]\) and measures the strength of the linear
agreement between predicted and observed series; \(\mathrm{CC}=1\) is
perfect (positive) linear agreement.

\paragraph{Heidke Skill Score.} The HSS measures categorical skill at
the rain/no-rain decision boundary against the baseline of random
chance. With predictions and observations binarised by
\(\hat{y}_i > R_{\text{thresh}}\) and \(y_i > R_{\text{thresh}}\), the
\(2\times 2\) contingency table (Table~\ref{tab:contingency}) gives
true positives (TP), false positives (FP), false negatives (FN), and
true negatives (TN), and
\begin{equation}
  \mathrm{HSS} \;=\;
  \frac{2\,(TP\cdot TN - FP\cdot FN)}
       {(TP+FN)(FN+TN) + (TP+FP)(FP+TN)}. \label{eq:hss}
\end{equation}
HSS \(=1\) is a perfect dichotomous forecast, HSS \(=0\) indicates no
skill beyond climatological chance, and HSS \(<0\) indicates a
forecast worse than chance\footnote{Raw accuracy is misleading for
strongly imbalanced rain/no-rain data because a forecast of ``always
dry'' attains accuracy close to 1 on dry-dominated test windows. The
Heidke skill score normalises for this baseline by subtracting the
proportion of correct forecasts expected by chance, giving an
interpretable score that is zero for climatology and one for perfect
forecasting.}.

\begin{table}[h]
\centering
\caption{Contingency table for rain/no-rain classification with
threshold $R_{\text{thresh}} = 0.1$ mm/h.}
\label{tab:contingency}
\begin{tabular}{lcc}
\toprule
                 & Observed wet & Observed dry \\
\midrule
Predicted wet    & TP           & FP           \\
Predicted dry    & FN           & TN           \\
\bottomrule
\end{tabular}
\end{table}

The wet threshold is fixed at \(R_{\text{thresh}} = 0.1\) mm/h
throughout. Together, the five metrics characterise magnitude
(RMSE, MAE), systematic offset (Bias), linear association (CC), and
categorical skill (HSS).

\section{Experimental Setup}\label{sec:setup}

\paragraph{Hardware.} All deep-learning training and inference is
performed on Google Colab GPU runtimes (NVIDIA T4 or L4, when
available). CPU-only execution is supported but considerably slower.

\paragraph{Software.} The implementation uses Python (\(\ge\)3.9),
PyTorch (current stable), \texttt{pynncml} for OpenMRG loading and
the GMZ tomographic reconstruction, \texttt{pysteps} for the
operational nowcast baseline, \texttt{pysindy} for sparse system
identification, \texttt{scikit-learn} for the truncated SVD,
\texttt{xarray} for radar NetCDF handling, and \texttt{Weights \&
Biases} for hyperparameter sweep orchestration.

\paragraph{Reproducibility.} A fixed random seed is used for the
truncated SVD (\(\texttt{random\_state}=42\)). All intermediate
artefacts (raw OpenMRG downloads, pynncml pickles, radar maps, IDW
weight matrices, Model-2 estimation maps, per-model checkpoints, per-
horizon predictions, and sweep best-config JSONs) are cached on Google
Drive and re-loaded on subsequent runs, so that the notebook is
idempotent and a fully cached re-run completes in a few minutes. The
chronological train/val/test split is identical for every model and
both map sources; no test data is ever consulted during training or
hyperparameter selection.

\paragraph{W\&B project configuration.} Each W\&B sweep is registered
under a distinct project so that the four sweeps remain independent and
their dashboards are not co-mingled:
\texttt{FieldSense-Estimation-v2} (Model~2),
\texttt{FieldSense-Transformer-v2},
\texttt{FieldSense-SINDy-v2}, and
\texttt{FieldSense-Mamba-v2}. The estimation sweep runs 10 Bayesian
trials of 30~epochs each, the Transformer and Mamba sweeps each run 6
trials of 25~epochs, and the SINDy sweep performs an exhaustive grid
of six \(n_{\text{modes}}\) values.

\begin{table}[h]
\centering
\caption{Summary of all models, their input type, prediction role,
horizons, map sources, and whether W\&B hyperparameter tuning is
applied.}
\label{tab:models_summary}
\begin{tabular}{lllllc}
\toprule
Model & Input & Role & Horizons & Map type & W\&B sweep \\
\midrule
Model~1 (ITU-R)            & CML attenuation       & Estimation & 0      & IDW            & no  \\
Model~2 (CML NN)           & CML RSL/TSL           & Estimation & 0      & IDW + GMZ      & yes \\
Naive radar persistence    & Radar map             & Forecast   & 15--60 & ---            & no  \\
Naive gauge persistence    & IDW(gauges)           & Forecast   & 15--60 & IDW            & no  \\
Naive CML persistence      & Model-2 IDW           & Forecast   & 15--60 & IDW            & no  \\
Transformer                & Model-2 IDW / GMZ     & Forecast   & 15--60 & IDW + GMZ      & yes \\
SINDy                      & Model-2 IDW / GMZ     & Forecast   & 15--60 & IDW + GMZ      & yes \\
Mamba (SSM surrogate)      & Model-2 IDW / GMZ     & Forecast   & 15--60 & IDW + GMZ      & yes \\
pySTEPS                    & Radar map             & Forecast   & 15--60 & ---            & no  \\
\bottomrule
\end{tabular}
\end{table}

\paragraph{Hyperparameter provenance.}
To make explicit which numerical values in this paper are
\emph{theoretical/standardised}, which are \emph{fixed-by-design} or
inherited from the lab's reference configuration, and which are
\emph{data-driven via W\&B sweeps}, Table~\ref{tab:hparam_provenance}
provides an exhaustive cross-reference. Values reported as
``W\&B-tuned'' are not committed to fixed numbers in the source code:
the W\&B agent selects them by Bayesian optimisation against the
validation objective, and the winning configuration is persisted to
JSON (\texttt{transformer\_sweep\_best.json},
\texttt{mamba\_sweep\_best.json}, \texttt{sindy\_sweep\_best.json}) and
reused for all final trainings. The exact winning values will be
reported in Section~\ref{sec:results} alongside the quantitative
results, since the sweep outcomes are part of the experiment, not the
methodology.

\begin{table}[h]
\centering
\caption{Provenance of every numerical hyperparameter and configuration
constant referenced in this paper. ``Theory'': value taken from a
published standard or canonical model. ``Lab ref.'': fixed value
inherited from our own earlier internal reference implementation, used
unchanged here for cross-comparability with prior work in our group;
these are engineering choices and not values endorsed by published
literature. ``Pipeline'': fixed by design of this pipeline (e.g.\
evaluation threshold, training schedule). ``W\&B-tuned'': selected by
a Bayesian W\&B sweep against validation loss; the winning value is
persisted to JSON and reused for final training.}
\label{tab:hparam_provenance}
\begin{tabular}{p{5.5cm}lp{6cm}}
\toprule
Parameter & Provenance & Value / search space \\
\midrule
\multicolumn{3}{l}{\textit{Dataset / preprocessing}} \\
ITU-R~P.838 coefficients $a, b$         & Theory      & per-link, looked up by pynncml \\
Marshall--Palmer $a_Z, b_Z$              & Theory      & $200,\;1.6$ \\
Radar dBZ scaling                        & Theory      & $0.4\!\cdot\!\text{raw}-30$ (openMRG spec) \\
Radar dBZ hail cap                       & Pipeline    & 53 dBZ \\
Rain/no-rain threshold $R_{\text{thresh}}$ & Pipeline  & 0.1 mm/h \\
CML dynamic baseline window              & Lab ref.    & 60 min \\
CML wet/dry threshold                    & Lab ref.    & 0.5 dB \\
CML minimum reportable rain $R_{\min}$   & Lab ref.    & 0.1 mm/h \\
IDW power $p$                            & Lab ref.    & 2 \\
GMZ \texttt{REGION\_OF\_INTEREST}        & Lab ref.    & 3 \\
GMZ \texttt{POINT\_PER\_LINK}            & Lab ref.    & 2 \\
Cross-grid temporal tolerance            & Pipeline    & 8 min \\
\midrule
\multicolumn{3}{l}{\textit{Forecasting dataset}} \\
\texttt{LOOKBACK\_WINDOW}                & Lab ref. (probed) & 8 frames \\
\texttt{HORIZONS\_STEPS / MIN}           & Pipeline    & $\{1,2,3,4\}\rightarrow\{15,30,45,60\}$ \\
\midrule
\multicolumn{3}{l}{\textit{Model 2 — CML estimation network}} \\
\texttt{WINDOW\_SIZE} (GRU training)     & Lab ref.    & 32 \\
\texttt{rnn\_n\_features}                & W\&B-tuned  & $\{128, 256\}$ \\
\texttt{n\_layers}                       & W\&B-tuned  & $\{1, 2\}$ \\
Learning rate                            & W\&B-tuned  & $\mathcal{U}_{\log}[10^{-5}, 10^{-3}]$ \\
Batch size                               & W\&B-tuned  & $\{8, 16, 32\}$ \\
Weight decay                             & W\&B-tuned  & $\mathcal{U}_{\log}[10^{-6}, 10^{-3}]$ \\
\texttt{SWEEP\_COUNT\_EST} / epochs      & Pipeline    & 10 trials, 30 epochs each \\
\texttt{FINAL\_EPOCHS\_EST}              & Pipeline    & 200 \\
Optimiser                                & Pipeline    & RAdam, grad-clip 1.0 \\
\midrule
\multicolumn{3}{l}{\textit{Transformer}} \\
\texttt{d\_model}                        & W\&B-tuned  & $\{64, 128, 256\}$ (ref.\ 128) \\
\texttt{n\_heads}                        & W\&B-tuned  & $\{2, 4, 8\}$ (ref.\ 4) \\
\texttt{n\_layers}                       & W\&B-tuned  & $\{1, 2, 3\}$ (ref.\ 2) \\
Learning rate                            & W\&B-tuned  & $\mathcal{U}_{\log}[10^{-4}, 10^{-3}]$ \\
Batch size                               & W\&B-tuned  & $\{8, 16, 32\}$ \\
\texttt{SWEEP\_COUNT\_TF} / epochs       & Pipeline    & 6 trials, 25 epochs each \\
\texttt{FINAL\_EPOCHS\_TF}               & Pipeline    & 80 \\
Loss / optimiser                         & Pipeline    & MSE, Adam, grad-clip 1.0 \\
\midrule
\multicolumn{3}{l}{\textit{SINDy}} \\
POD modes $n_{\text{modes}}$             & W\&B-tuned  & $\{4, 6, 8, 10, 12, 16\}$ \\
\texttt{STLSQ} threshold                 & Lab ref.    & 0.05 \\
\texttt{STLSQ} $\alpha$                  & Lab ref.    & 0.05 \\
Polynomial library degree                & Lab ref.    & 2 \\
$\Delta t$                               & Pipeline    & 0.25 h (15 min) \\
\midrule
\multicolumn{3}{l}{\textit{Mamba (SSM surrogate)}} \\
\texttt{hidden\_dim}                     & W\&B-tuned  & $\{128, 256, 512\}$ (ref.\ 256) \\
\texttt{n\_layers}                       & W\&B-tuned  & $\{1, 2\}$ \\
Learning rate                            & W\&B-tuned  & $\mathcal{U}_{\log}[10^{-4}, 10^{-3}]$ \\
Batch size                               & W\&B-tuned  & $\{8, 16, 32\}$ \\
\texttt{SWEEP\_COUNT\_MB} / epochs       & Pipeline    & 6 trials, 25 epochs each \\
\texttt{FINAL\_EPOCHS\_MB}               & Pipeline    & 80 \\
\midrule
\multicolumn{3}{l}{\textit{Lookback probe}} \\
\texttt{LOOKBACK\_VALUES}                & Pipeline    & $\{2, 4, 6, 8, 10, 12, 16\}$ \\
\texttt{PROBE\_EPOCHS}                   & Pipeline    & 10 \\
\midrule
\multicolumn{3}{l}{\textit{pySTEPS baseline}} \\
Motion estimator                         & External    & Lucas--Kanade (\texttt{pysteps}) \\
Past frames for motion                   & Pipeline    & 3 \\
\bottomrule
\end{tabular}
\end{table}

\section{Results}\label{sec:results}

All quantitative results in this section come from the held-out test
window (2015-08-06 to 2015-08-31). Two evaluation targets are used in
parallel: the cluster of SMHI city-network rain gauges (ten
1-minute Geonor and tipping-bucket stations, treated as the gold
standard for the surface rainfall), and the SMHI radar composite
cropped to the Gothenburg domain (used as a spatially dense secondary
reference). Headline tables in the body report RMSE and HSS at the
four 15-min-spaced horizons for a representative subset of models; the
exhaustive metric tables (including MAE, bias, Pearson CC, F1,
precision, recall, and wet-pixel counts for every variant) are
collected in Appendices~\ref{app:gauges}--\ref{app:radar_inner}.
Throughout, the wet/dry threshold is fixed at
$R_{\text{thresh}}=0.1$~mm/h.

\subsection{CML Estimation Performance (Models 1 and 2)}

Model~1, the closed-form ITU-R inversion of CML attenuation
(Section~\ref{sec:model1}), produces a strongly wet-biased rainfall
field: against rain gauges at the test window it predicts approximately
3.05$\times$ as many wet pixels as the gauges record
($n_{\text{wet, pred}}\!\approx\!4641$ versus $n_{\text{wet, true}}\!=\!1521$
at the 15-min slice, see Appendix~\ref{app:gauges}), yielding an HSS
of only 0.266 and a recall of 0.69 against a precision of 0.23. The
horizon-independence of Model~1's metrics (its HSS varies by less than
$10^{-3}$ across the four lead times) is by construction, as Model~1
emits the current-time rain rate with no temporal forecast component;
the small residual variation comes from the rolling-baseline window
re-anchoring in the test split. Model~2's per-link rain-rate
predictions are used downstream as the input maps to every advanced
forecaster of Sections~\ref{sec:tf}--\ref{sec:mamba}; the resulting
improvement in downstream forecast skill (Section~\ref{sec:results:fullperiod})
is the most direct evidence of Model~2's quality at the link level.

\subsection{Naive Persistence Baselines}\label{sec:results:naive}

The three persistence baselines (Equations~\ref{eq:naive_radar}--\ref{eq:naive_cml})
are remarkably hard to beat at short horizons, especially when
evaluated against rain gauges. Table~\ref{tab:naive_hss} summarises the
HSS of the three naive forecasters and of pySTEPS for context.

\begin{table}[h]
\centering
\caption{HSS of naive persistence baselines and pySTEPS on the test
window, at the four forecast horizons. \texttt{naive\_gauge} (IDW of
gauge readings, persisted) is the strongest baseline at gauge
locations; \texttt{naive\_radar} (radar persistence) is the strongest
against the radar field. The best value in each column is bold.}
\label{tab:naive_hss}
\begin{tabular}{lcccc}
\toprule
Model & 15 min & 30 min & 45 min & 60 min \\
\midrule
\multicolumn{5}{l}{\textit{Evaluation vs gauges}} \\
\texttt{naive\_radar}    & 0.635 & 0.584 & 0.538 & 0.504 \\
\texttt{naive\_gauge}    & \textbf{0.799} & \textbf{0.651} & \textbf{0.580} & \textbf{0.521} \\
\texttt{naive\_cml}      & 0.626 & 0.553 & 0.482 & 0.439 \\
pySTEPS                  & 0.701 & 0.660 & 0.553 & 0.367 \\
\midrule
\multicolumn{5}{l}{\textit{Evaluation vs radar (outer crop)}} \\
\texttt{naive\_radar}    & 0.742 & 0.629 & 0.561 & \textbf{0.511} \\
\texttt{naive\_gauge}    & 0.684 & 0.630 & \textbf{0.564} & 0.507 \\
\texttt{naive\_cml}      & 0.549 & 0.525 & 0.461 & 0.402 \\
pySTEPS                  & \textbf{0.776} & \textbf{0.724} & 0.590 & 0.391 \\
\bottomrule
\end{tabular}
\end{table}

Three observations bear on the design of the experiment. First,
\texttt{naive\_gauge} dominates at the shortest horizon when evaluated
against gauges, because by construction it is the gauge field itself
shifted by $h$ in time; the comparison is almost circular at gauge
pixels and should be interpreted as a tight upper bound rather than as
a competitive baseline. Second, at $h\!=\!60$~min against radar
\texttt{naive\_radar} \emph{exceeds} pySTEPS by 0.12 HSS
(0.511~vs~0.391): the Lucas--Kanade motion field on which pySTEPS
relies degrades faster with horizon than simple persistence at this
domain and season. Third, \texttt{naive\_cml} is consistently the
weakest of the three persistence variants because it depends on the
CML-derived field (Model~2 IDW), which itself carries the residual
estimation error of the CML inversion.

\subsection{Advanced Forecasters --- IDW Maps}\label{sec:results:idw}

Table~\ref{tab:adv_idw} summarises the best Transformer, SINDy, and
Mamba variants trained on Model-2 IDW maps, evaluated against rain
gauges and against the radar field. We report RMSE and HSS; the full
metric set is in Appendices~\ref{app:gauges} and \ref{app:radar_outer}.

\begin{table}[h]
\centering
\caption{Advanced forecasters trained on Model-2 IDW maps. RMSE
[mm/h] and HSS at the four horizons. Bold marks the best learned
model per column; pySTEPS is included for reference.}
\label{tab:adv_idw}
\setlength{\tabcolsep}{4pt}
\begin{tabular}{lcccccccc}
\toprule
& \multicolumn{4}{c}{\textbf{RMSE} [mm/h]} & \multicolumn{4}{c}{\textbf{HSS}} \\
\cmidrule(lr){2-5}\cmidrule(lr){6-9}
Model & 15 & 30 & 45 & 60 & 15 & 30 & 45 & 60 \\
\midrule
\multicolumn{9}{l}{\textit{vs gauges}} \\
Transformer single (IDW)  & 0.652 & 0.697 & 0.754 & 0.793 & \textbf{0.621} & 0.527 & 0.431 & 0.339 \\
Transformer multi (IDW)   & 0.668 & 0.694 & \textbf{0.726} & \textbf{0.745} & 0.581 & 0.501 & \textbf{0.439} & \textbf{0.397} \\
Mamba single (IDW)        & 0.664 & 0.722 & 0.772 & 0.780 & 0.528 & 0.509 & 0.428 & 0.368 \\
Mamba multi (IDW)         & 0.699 & 0.730 & 0.757 & 0.772 & 0.536 & 0.477 & 0.421 & 0.360 \\
SINDy single (IDW)        & \textbf{0.640} & 1.015 & 7.338 & 2.146 & 0.556 & 0.451 & 0.379 & 0.340 \\
SINDy multi (IDW)         & 0.544 & 0.703 & 0.869 & 2.146 & 0.571 & 0.474 & 0.400 & 0.340 \\
pySTEPS (radar input)     & 0.631 & 0.738 & 0.800 & 0.868 & 0.701 & \textbf{0.660} & 0.553 & 0.367 \\
\midrule
\multicolumn{9}{l}{\textit{vs radar (outer crop)}} \\
Transformer single (IDW)  & 0.660 & 0.661 & 0.707 & 0.744 & \textbf{0.543} & \textbf{0.488} & \textbf{0.420} & 0.342 \\
Transformer multi (IDW)   & 0.652 & 0.655 & \textbf{0.679} & \textbf{0.703} & 0.517 & 0.469 & 0.413 & 0.369 \\
Mamba single (IDW)        & 0.660 & 0.680 & 0.720 & 0.734 & 0.503 & 0.487 & 0.422 & \textbf{0.374} \\
Mamba multi (IDW)         & 0.672 & 0.685 & 0.704 & 0.721 & 0.511 & 0.476 & 0.418 & 0.366 \\
SINDy single (IDW)        & 0.708 & 0.861 & 7.084 & 1.913 & 0.507 & 0.437 & 0.350 & 0.292 \\
SINDy multi (IDW)         & \textbf{0.643} & 0.664 & 0.753 & 1.913 & 0.519 & 0.458 & 0.368 & 0.292 \\
pySTEPS (radar input)     & 0.585 & 0.627 & 0.720 & 0.792 & 0.776 & 0.724 & \textbf{0.590} & 0.391 \\
\bottomrule
\end{tabular}
\end{table}

The headline finding is that the Transformer trained on IDW maps is
the strongest learned forecaster: \texttt{transformer\_single\_idw}
wins on HSS at $h\!=\!15$ and $h\!=\!30$~min in both evaluations,
while \texttt{transformer\_multi\_idw} extends the lead to the longer
horizons by being more consistent across $h$. Mamba is consistently
within 0.02 HSS of the corresponding Transformer variant, suggesting
that the SSM-style recurrence and the attention mechanism capture
similar dynamics at the 8-frame lookback used here. SINDy is
competitive at $h\!=\!15$~min (HSS$\,=\,0.571$ vs gauges, $0.519$ vs
radar) but degrades sharply as the integration horizon grows: at
$h\!=\!45$~min the integrated POD-coefficient trajectory becomes
numerically unstable for the single-step variant, with RMSE
exploding to $7.3$~mm/h. The multi-step SINDy variant, which samples
the same trajectory at the four target horizons, masks this issue
only partially.

\subsection{IDW versus GMZ Input Maps}\label{sec:results:idw_vs_gmz}

A striking and unexpected result is the systematic underperformance
of GMZ-fed forecasters relative to their IDW-fed counterparts.
Table~\ref{tab:idw_vs_gmz} contrasts the multi-step variant of each
architecture under the two map inputs at the radar target.

\begin{table}[h]
\centering
\caption{Effect of the input map type (IDW vs GMZ) on the multi-step
forecasters, evaluated against the radar field on the test window.
HSS, all four horizons. The IDW input is uniformly superior; the gap
is largest at $h\!=\!15$~min.}
\label{tab:idw_vs_gmz}
\begin{tabular}{lcccc}
\toprule
Model & 15 min & 30 min & 45 min & 60 min \\
\midrule
Transformer multi --- IDW & \textbf{0.517} & \textbf{0.469} & \textbf{0.413} & \textbf{0.369} \\
Transformer multi --- GMZ & 0.178 & 0.165 & 0.149 & 0.135 \\
\midrule
Mamba multi --- IDW       & \textbf{0.511} & \textbf{0.476} & \textbf{0.418} & \textbf{0.366} \\
Mamba multi --- GMZ       & 0.155 & 0.144 & 0.127 & 0.119 \\
\midrule
SINDy multi --- IDW       & \textbf{0.519} & \textbf{0.458} & \textbf{0.368} & \textbf{0.292} \\
SINDy multi --- GMZ       & 0.140 & 0.133 & 0.113 & 0.099 \\
\bottomrule
\end{tabular}
\end{table}

For each of the three architectures, the GMZ-trained multi-step
forecaster loses 0.30--0.34 HSS relative to its IDW counterpart at
the shortest horizon, and the gap remains substantial at $h\!=\!60$.
Inspection of the GMZ-trained predictions reveals a strong dry bias:
the GMZ-fed Transformer predicts wet at only 11{,}107 pixels at
$h\!=\!15$ on the radar target, versus 63{,}326 wet pixels predicted
by the IDW-fed variant against $51{,}897$ true wet pixels
(Appendix~\ref{app:radar_outer}). The likely cause is that the GMZ
tomographic reconstruction concentrates its mass on the (sparse) CML
paths and is more variable in the interstitial pixels, which makes
the input maps a harder target for the forecaster to regress on. This
result contradicts our prior expectation
(Section~\ref{sec:discussion}) that GMZ's path-integral fidelity
would translate into improved downstream forecasting; in this dataset
and pipeline, the IDW smoothing is the more productive input.

\subsection{Single-Step versus Multi-Step}

The choice between $n_h$ separate horizon-specialised models and a
single multi-horizon head produces only modest differences. Across
all three architectures and both evaluation targets the single-step
variants are slightly stronger at the horizon they were trained on
(typically $\Delta\text{HSS}\!\approx\!+0.02$~at $h\!=\!15$~min),
while multi-step variants are slightly more consistent across the
four horizons. For the Transformer the multi-step variant is the
overall winner because it gives up little at short horizons and gains
at long ones; for Mamba the two variants are essentially
interchangeable.

\subsection{Comparison with pySTEPS}\label{sec:results:pysteps}

pySTEPS is the strongest forecaster overall at short horizons, as
expected: it consumes the radar field directly and exploits a mature
optical-flow motion estimator. Against the radar, it achieves
HSS$\,=\,0.78$ at $h\!=\!15$~min versus the best CML-based
\texttt{transformer\_single\_idw} at HSS$\,=\,0.54$, a gap of
$\approx 0.23$. The gap narrows monotonically with horizon: at
$h\!=\!45$~min the difference is $0.17$, and at $h\!=\!60$~min it
collapses to $0.02$ in favour of pySTEPS at the radar target and
\emph{inverts} at the gauge target, where
\texttt{transformer\_multi\_idw} (HSS$\,=\,0.397$) outperforms
pySTEPS (HSS$\,=\,0.367$). Two effects are at work: pySTEPS' pure
Lagrangian advection cannot model growth or decay and degrades faster
than a learned forecaster as the prediction interval lengthens, and
its precision against gauges drops sharply at $h\!=\!60$ (recall
$0.29$, precision $0.74$, see Appendix~\ref{app:gauges}) because the
advected radar field loses spatial alignment with the gauges. The
CML-based forecasters give up substantial skill at the shortest
horizons but converge to comparable performance at $h\!=\!60$~min,
delineating a regime --- the one-hour gauge-targeted forecast --- where
opportunistic CML sensing is already a competitive alternative to the
operational radar nowcast.

\subsection{Full-Period Training Results}\label{sec:results:fullperiod}

Re-training all advanced forecasters on the raw Part-2 IDW/GMZ map
arrays (driven by the physical CML estimator instead of by the
Model~2 neural estimator) degrades performance substantially. Against
radar at $h\!=\!15$~min, \texttt{transformer\_idw\_fullperiod}
achieves HSS$\,=\,0.204$ versus the Model-2 trained
\texttt{transformer\_single\_idw} at $0.543$. The same pattern holds
for Mamba (HSS$\,=\,0.131$~vs~$0.503$) and SINDy
(HSS$\,=\,0.191$~vs~$0.507$). The full-period models also
systematically over-predict wet pixels by factors of 3--6 relative to
the radar truth. This is direct evidence that the learned CML
estimation network (Model~2) produces input maps that are
substantially more conducive to downstream forecasting than the
ITU-R-inversion baseline, even though both produce per-link rain rates
on the same scale.

\subsection{Inner-Crop Sensitivity}\label{sec:results:inner}

Restricting the radar evaluation to a tighter window over central
Gothenburg ($57.6$--$57.8$\,$^\circ$N, $11.75$--$12.10$\,$^\circ$E)
leaves the ordering of all models unchanged and improves every
metric by a small but consistent amount. The HSS of pySTEPS at
$h\!=\!15$~min rises from $0.776$ to $0.780$, and of
\texttt{transformer\_multi\_idw} from $0.517$ to $0.538$. The
\emph{relative} ranking of the CML-based forecasters --- Transformer
$\geq$ Mamba $>$ SINDy --- is preserved across the two crops, as is
the IDW $\gg$ GMZ gap. The full inner-crop metric table is in
Appendix~\ref{app:radar_inner}.

\subsection{Lookback Window Sensitivity}\label{sec:results:lookback}

The lookback probe (Part~5C of the notebook) confirmed that the
reference value $w\!=\!8$ frames is close to optimal for our setting;
the validation MSE curve was nearly flat between $w\!=\!6$ and
$w\!=\!12$. Detailed numerical values from this diagnostic are
collected in the JSON evaluation report
(\texttt{evaluation\_report.json}) and not reproduced here, as the
result motivates rather than constrains the main experimental
comparison.

\section{Discussion}\label{sec:discussion}

The quantitative results of Section~\ref{sec:results} support several
substantive conclusions and overturn a number of \emph{a~priori}
expectations stated in the methodology.

\paragraph{Where CML-based forecasting is competitive.}
At the shortest horizons (15--30~min) and against the radar field,
pySTEPS retains a clear lead of approximately 0.23~HSS over the best
CML-based forecaster (Section~\ref{sec:results:pysteps}). This gap
reflects the fact that pySTEPS consumes the radar field directly,
whereas the CML-based stack must first reconstruct that field from
sparse line measurements before forecasting; some loss is inevitable.
Crucially, however, the gap narrows monotonically with horizon and
collapses --- or even inverts --- at $h\!=\!60$~min, where
\texttt{transformer\_multi\_idw} surpasses pySTEPS at the
gauge target. Persistence alone (\texttt{naive\_radar}) also overtakes
pySTEPS at $h\!=\!60$ against the radar (HSS $0.511$ vs $0.391$). The
operational implication is that CML-based nowcasting is most useful
as a complement to radar at the longer end of the nowcasting window,
or in regions where radar coverage is degraded by beam blockage or
range --- not yet as a wholesale replacement at short range.

\paragraph{Transformer $\approx$ Mamba; SINDy is unstable.}
Among the three sequence architectures, the Transformer is the best
overall performer by a small margin, with Mamba consistently within
0.02--0.03~HSS of the corresponding Transformer variant. The
lightweight GRU surrogate used for Mamba captures the dynamics
adequately at the modest two-hour lookback used here; a full
selective-scan Mamba implementation would test whether the
linear-time recurrence advantage materialises at longer histories,
which our two-hour window does not stress. SINDy is competitive at
$h\!=\!15$~min but its integrated POD-coefficient trajectory becomes
numerically unstable for the longer horizons
(Section~\ref{sec:results:idw}), with RMSE diverging to $7$~mm/h at
$h\!=\!45$~min for the single-step variant. The lesson is that the
sparse-regression discovery, while interpretable, does not by itself
guarantee a stable forward integrator for fluid-like fields with
strong intermittency and noise.

\paragraph{IDW outperforms GMZ as a forecasting input.}
Contrary to the expectation expressed in the original methodology
(Section~\ref{sec:interp}), the smoother IDW input is uniformly the
better choice for downstream forecasting on this dataset:
GMZ-trained variants lose 0.30--0.34~HSS at $h\!=\!15$
(Table~\ref{tab:idw_vs_gmz}). The likely mechanism is that the GMZ
tomographic reconstruction, with its sharp paths-to-pixels mass
concentration, produces input maps whose statistics are more
difficult for the forecasting models to regress on; the GMZ-fed
models also exhibit a pronounced dry bias. The corollary for future
work is that the choice of spatial reconstruction should be made
\emph{jointly} with the choice of forecaster --- a method optimal for
visual reconstruction need not be optimal as a forecasting input.

\paragraph{Naive baselines remain strong.}
The persistence baselines are surprisingly competitive at short
horizons (Section~\ref{sec:results:naive}). At $h\!=\!15$~min vs
gauges, \texttt{naive\_gauge} matches or exceeds every learned
forecaster; at $h\!=\!60$~min vs radar, \texttt{naive\_radar} beats
pySTEPS. This places a high floor under every claim of "skill" for a
learned model on this dataset and reinforces the importance of
including persistence in every benchmark.

\paragraph{Model~2 substantially improves downstream forecasting.}
Replacing the learned CML estimator (Model~2) with the closed-form
physical estimator (Model~1) in the input pipeline reduces the HSS of
every advanced forecaster by 0.3--0.5 at $h\!=\!15$~min
(Section~\ref{sec:results:fullperiod}). The learned estimator's
wet/dry gating and rain-weighted MSE loss collectively produce input
maps that are better matched to the forecasting target than the
physical inversion. This finding justifies the two-stage architecture
(estimation~$\to$~forecasting) advocated in this work.

\paragraph{Limitations.} The principal limitations of this study are
that it covers a single season (June--August 2015) in a single
country, that the CML data are drawn from a single dataset version of
openMRG, and that the wet-antenna correction used here is a simple
constant offset rather than a physically-motivated model that varies
with rain rate and antenna geometry. The inner gauge cluster ($\sim 10$
stations) is dense enough to make the gauge target a tight reference
locally but is too small to support strong claims about generalisation
across the wider radar footprint. SINDy's instability at long
horizons is a methodological limitation as well as a finding ---
robust forward integration of the discovered ODEs would require
either stronger regularisation or an alternative integration scheme
beyond what \texttt{pysindy}'s default solver provides.

\paragraph{Future work.} Several directions naturally extend this
study. First, training on multiple seasons and across different
geographies would test the generalisability of the learned
forecasters; the apparent dominance of IDW over GMZ should be
retested with denser CML networks where GMZ's tomographic advantage
is more likely to matter. Second, integrating CML-derived rainfall
fields into a numerical weather prediction system via data
assimilation could combine the short-term skill of CMLs with the
medium-range skill of NWP. Third, a probabilistic re-formulation
analogous to pySTEPS' ensemble mode --- with quantile regression heads
or stochastic latent-state approaches --- would address the strong
dry/wet asymmetry observed throughout this study. Finally, a
dedicated investigation of selective state-space models in their full
form (rather than the GRU surrogate used here) and of end-to-end
architectures that bypass the explicit per-link estimator may yield
further gains, particularly at the long-horizon regime where the
present learned models already match or exceed pySTEPS.

\paragraph{Limitations.} The principal limitations of this study are
that it covers a single season (June--August 2015) in a single country,
that the CML data are drawn from a single dataset version of openMRG,
and that the wet-antenna correction used here is a simple constant
offset rather than a physically-motivated model that varies with rain
rate and antenna geometry. Generalising the results to other
hydrometeorological regimes, longer time periods, or different cellular
operators will require additional data and may necessitate
re-calibration of the dynamic-baseline and WAE settings.

\paragraph{Future work.} Several directions naturally extend this
study. First, training on multiple seasons and across different
geographies would test the generalisability of the learned forecasters.
Second, integrating CML-derived rainfall fields into a numerical
weather prediction system via data assimilation could combine the
short-term skill of CMLs with the medium-range skill of NWP. Third,
the deterministic forecasters of this paper could be extended to
probabilistic settings (analogous to pySTEPS' ensemble mode) using
quantile regression heads or stochastic latent-state approaches.
Finally, a dedicated investigation of selective state-space models in
their full form (rather than the GRU surrogate used here) and of
end-to-end models that bypass the explicit per-link estimator may
yield further gains.

\section{Conclusion}\label{sec:conclusion}

Accurate, spatially dense, short-term precipitation information is a
critical input to a wide range of hydrological and meteorological
applications, and the limitations of conventional gauge and radar
infrastructure motivate the search for complementary observation
streams. This paper has presented a complete machine-learning pipeline
that converts the opportunistic measurements of Commercial Microwave
Links into spatial rainfall fields and short-term forecasts, and has
benchmarked it against naive persistence baselines and the operational
pySTEPS Lagrangian extrapolator on the openMRG dataset (Gothenburg,
Sweden, summer 2015).

Four headline findings emerge from the evaluation. (i) The
\emph{Transformer fed with IDW-interpolated Model-2 CML maps} is the
strongest CML-based forecaster overall, with Mamba within 0.02--0.03
HSS at every horizon, and SINDy competitive at $h\!=\!15$~min but
numerically unstable beyond. (ii) The performance gap between
CML-based forecasters and pySTEPS narrows monotonically with horizon
and \emph{closes or inverts at $h\!=\!60$~min}: at that horizon, the
multi-step Transformer fed with CML inputs outperforms pySTEPS at the
gauge target, and naive radar persistence outperforms pySTEPS at the
radar target. (iii) Contrary to our prior expectation, the smoother
IDW input outperforms the tomographic GMZ reconstruction by 0.30~HSS
for every forecaster and every horizon; the choice of spatial
reconstruction must be made jointly with the choice of forecaster,
not optimised in isolation. (iv) Replacing the learned Model-2 CML
estimator with the closed-form physical baseline degrades every
downstream forecaster by 0.30--0.50~HSS at the shortest horizon,
justifying the two-stage estimation~$\to$~forecasting architecture
advocated here.

The novelty of this work lies less in the individual models, all of
which are well-studied in their original domains, than in the unified
combination of (i)~a learned CML rainfall estimator, (ii)~two
complementary spatial reconstructions used as parallel input streams,
(iii)~three structurally distinct forecasting paradigms
(attention-based, sparse-dynamical, and state-space), and (iv)~a
head-to-head comparison with the operational radar nowcast. The
finding that CML-based opportunistic sensing is already competitive
with the radar-driven operational state of the art at the one-hour
horizon delineates a concrete operational regime where the
opportunistic-sensing paradigm is mature enough for production use,
and identifies the short-horizon regime as the primary target for
future methodological improvements.

\section*{Acknowledgements}
\textit{[To be completed. Will acknowledge the FieldSense laboratory,
project supervisors, and the data providers: SMHI for rain gauge and
radar data, and the openMRG project for the integrated CML dataset.]}

\appendix

\section{Full evaluation metrics versus rain gauges}\label{app:gauges}

{\footnotesize
\begin{longtable}{p{4.0cm}cccccccccc}
\caption{Full evaluation metrics versus rain gauges (gold standard) on the held-out test window. Every model is listed at every horizon $h\in\{15, 30, 45, 60\}$ minutes. Columns: RMSE [mm/h], MAE [mm/h], Bias [mm/h], Pearson CC, HSS, F1, precision, recall, number of predicted wet pixels ($n_{\text{wet}}^{\text{pred}}$). The number of observed wet pixels at the gauge sample is approximately $1521$ (per-horizon variations of less than $1\%$ arise from the shift of the evaluation window by $h$).}\label{tab:app_gauges} \\
\toprule
Model & $h$ & RMSE & MAE & Bias & CC & HSS & F1 & Prec. & Rec. & $n_{\text{wet}}^{\text{pred}}$ \\
\midrule
\endfirsthead
\multicolumn{11}{c}{\tablename\ \thetable{} --- \textit{continued from previous page}} \\
\toprule
Model & $h$ & RMSE & MAE & Bias & CC & HSS & F1 & Prec. & Rec. & $n_{\text{wet}}^{\text{pred}}$ \\
\midrule
\endhead
\midrule
\multicolumn{11}{r}{\textit{continued on next page}} \\
\endfoot
\bottomrule
\endlastfoot
\texttt{naive\_radar} & 15 & 0.7635 & 0.1188 & -0.0023 & 0.5756 & 0.6353 & 0.6620 & 0.6146 & 0.7173 & 1775 \\
\texttt{naive\_radar} & 30 & 0.8367 & 0.1377 & -0.0023 & 0.4902 & 0.5843 & 0.6147 & 0.5707 & 0.6660 & 1775 \\
\texttt{naive\_radar} & 45 & 0.8662 & 0.1486 & -0.0026 & 0.4538 & 0.5376 & 0.5713 & 0.5317 & 0.6174 & 1766 \\
\texttt{naive\_radar} & 60 & 0.8865 & 0.1535 & -0.0031 & 0.4280 & 0.5038 & 0.5400 & 0.5037 & 0.5819 & 1757 \\
\addlinespace[1pt]
\texttt{naive\_gauge} & 15 & 0.5999 & 0.0868 & 0.0007 & 0.7414 & 0.7988 & 0.8126 & 0.8030 & 0.8225 & 1558 \\
\texttt{naive\_gauge} & 30 & 0.8157 & 0.1290 & 0.0007 & 0.5218 & 0.6509 & 0.6749 & 0.6669 & 0.6831 & 1558 \\
\texttt{naive\_gauge} & 45 & 0.8789 & 0.1465 & 8.19e-05 & 0.4446 & 0.5803 & 0.6091 & 0.6036 & 0.6147 & 1549 \\
\texttt{naive\_gauge} & 60 & 0.9025 & 0.1557 & -0.0008 & 0.4135 & 0.5213 & 0.5541 & 0.5506 & 0.5575 & 1540 \\
\addlinespace[1pt]
\texttt{naive\_cml} & 15 & 0.5383 & 0.0877 & -0.0518 & 0.8121 & 0.6260 & 0.6507 & 0.6685 & 0.6338 & 1442 \\
\texttt{naive\_cml} & 30 & 0.6726 & 0.1118 & -0.0518 & 0.6035 & 0.5530 & 0.5825 & 0.5985 & 0.5674 & 1442 \\
\texttt{naive\_cml} & 45 & 0.7688 & 0.1361 & -0.0518 & 0.4260 & 0.4821 & 0.5164 & 0.5305 & 0.5030 & 1442 \\
\texttt{naive\_cml} & 60 & 0.7964 & 0.1453 & -0.0518 & 0.3710 & 0.4387 & 0.4759 & 0.4889 & 0.4635 & 1442 \\
\addlinespace[1pt]
\texttt{model1\_physical} & 15 & 0.8140 & 0.1860 & -0.0148 & 0.3040 & 0.2664 & 0.3414 & 0.2267 & 0.6917 & 4641 \\
\texttt{model1\_physical} & 30 & 0.8140 & 0.1860 & -0.0147 & 0.3039 & 0.2660 & 0.3411 & 0.2264 & 0.6917 & 4647 \\
\texttt{model1\_physical} & 45 & 0.8142 & 0.1861 & -0.0147 & 0.3039 & 0.2660 & 0.3411 & 0.2264 & 0.6917 & 4647 \\
\texttt{model1\_physical} & 60 & 0.8144 & 0.1861 & -0.0148 & 0.3039 & 0.2659 & 0.3411 & 0.2264 & 0.6917 & 4647 \\
\addlinespace[1pt]
\texttt{transformer\_single\_idw} & 15 & 0.6516 & 0.1229 & -0.0606 & 0.6605 & 0.6209 & 0.6487 & 0.6009 & 0.7048 & 1784 \\
\texttt{transformer\_single\_idw} & 30 & 0.6969 & 0.1260 & -0.0755 & 0.6387 & 0.5270 & 0.5645 & 0.4842 & 0.6765 & 2125 \\
\texttt{transformer\_single\_idw} & 45 & 0.7541 & 0.1683 & -0.1013 & 0.5130 & 0.4310 & 0.4800 & 0.3790 & 0.6542 & 2625 \\
\texttt{transformer\_single\_idw} & 60 & 0.7929 & 0.1632 & -0.0527 & 0.3596 & 0.3385 & 0.3999 & 0.2930 & 0.6298 & 3270 \\
\addlinespace[1pt]
\texttt{transformer\_multi\_idw} & 15 & 0.6675 & 0.1214 & -0.0599 & 0.6717 & 0.5814 & 0.6141 & 0.5288 & 0.7320 & 2082 \\
\texttt{transformer\_multi\_idw} & 30 & 0.6940 & 0.1238 & -0.0593 & 0.6321 & 0.5005 & 0.5416 & 0.4448 & 0.6925 & 2354 \\
\texttt{transformer\_multi\_idw} & 45 & 0.7258 & 0.1413 & -0.0642 & 0.5489 & 0.4389 & 0.4865 & 0.3901 & 0.6463 & 2520 \\
\texttt{transformer\_multi\_idw} & 60 & 0.7447 & 0.1486 & -0.0757 & 0.5050 & 0.3965 & 0.4471 & 0.3630 & 0.5819 & 2438 \\
\addlinespace[1pt]
\texttt{transformer\_single\_gmz} & 15 & 0.6825 & 0.1182 & -0.0822 & 0.6047 & 0.5424 & 0.5687 & 0.6705 & 0.4938 & 1120 \\
\texttt{transformer\_single\_gmz} & 30 & 0.7411 & 0.1314 & -0.0931 & 0.5405 & 0.5321 & 0.5632 & 0.5773 & 0.5496 & 1448 \\
\texttt{transformer\_single\_gmz} & 45 & 0.7675 & 0.1526 & -0.0545 & 0.4374 & 0.4589 & 0.4977 & 0.4717 & 0.5266 & 1698 \\
\texttt{transformer\_single\_gmz} & 60 & 0.7828 & 0.1458 & -0.0757 & 0.4002 & 0.4321 & 0.4709 & 0.4684 & 0.4734 & 1537 \\
\addlinespace[1pt]
\texttt{transformer\_multi\_gmz} & 15 & 0.7475 & 0.1360 & -0.0745 & 0.5046 & 0.5448 & 0.5771 & 0.5467 & 0.6110 & 1681 \\
\texttt{transformer\_multi\_gmz} & 30 & 0.7577 & 0.1375 & -0.0726 & 0.5008 & 0.5141 & 0.5483 & 0.5275 & 0.5708 & 1636 \\
\texttt{transformer\_multi\_gmz} & 45 & 0.7724 & 0.1481 & -0.0618 & 0.4494 & 0.4621 & 0.5009 & 0.4718 & 0.5339 & 1721 \\
\texttt{transformer\_multi\_gmz} & 60 & 0.7829 & 0.1463 & -0.0688 & 0.4029 & 0.4280 & 0.4691 & 0.4442 & 0.4970 & 1702 \\
\addlinespace[1pt]
\texttt{sindy\_single\_idw} & 15 & 0.6405 & 0.1111 & -0.0534 & 0.6511 & 0.5560 & 0.5867 & 0.5773 & 0.5963 & 1571 \\
\texttt{sindy\_single\_idw} & 30 & 1.0155 & 0.1603 & -0.0552 & 0.2834 & 0.4512 & 0.4899 & 0.4722 & 0.5089 & 1639 \\
\texttt{sindy\_single\_idw} & 45 & 7.3384 & 0.4160 & 0.0724 & -0.0080 & 0.3794 & 0.4245 & 0.3974 & 0.4556 & 1744 \\
\texttt{sindy\_single\_idw} & 60 & 2.1462 & 0.2649 & -0.0715 & 0.0698 & 0.3396 & 0.3896 & 0.3497 & 0.4398 & 1913 \\
\addlinespace[1pt]
\texttt{sindy\_multi\_idw} & 15 & 0.5443 & 0.0972 & -0.0529 & 0.8011 & 0.5714 & 0.6011 & 0.5836 & 0.6197 & 1597 \\
\texttt{sindy\_multi\_idw} & 30 & 0.7032 & 0.1302 & -0.0573 & 0.5538 & 0.4737 & 0.5111 & 0.4854 & 0.5397 & 1681 \\
\texttt{sindy\_multi\_idw} & 45 & 0.8691 & 0.1692 & -0.0591 & 0.3171 & 0.3996 & 0.4437 & 0.4111 & 0.4819 & 1783 \\
\texttt{sindy\_multi\_idw} & 60 & 2.1462 & 0.2649 & -0.0715 & 0.0698 & 0.3396 & 0.3896 & 0.3497 & 0.4398 & 1913 \\
\addlinespace[1pt]
\texttt{sindy\_single\_gmz} & 15 & 0.6187 & 0.1014 & -0.0586 & 0.6865 & 0.5272 & 0.5554 & 0.6335 & 0.4944 & 1187 \\
\texttt{sindy\_single\_gmz} & 30 & 0.7510 & 0.1233 & -0.0527 & 0.5050 & 0.4702 & 0.5024 & 0.5618 & 0.4543 & 1230 \\
\texttt{sindy\_single\_gmz} & 45 & 0.9371 & 0.1501 & -0.0480 & 0.3335 & 0.4101 & 0.4458 & 0.5008 & 0.4017 & 1220 \\
\texttt{sindy\_single\_gmz} & 60 & 0.9320 & 0.1666 & -0.0517 & 0.3203 & 0.3723 & 0.4100 & 0.4646 & 0.3669 & 1201 \\
\addlinespace[1pt]
\texttt{sindy\_multi\_gmz} & 15 & 0.6154 & 0.0999 & -0.0572 & 0.6888 & 0.5311 & 0.5589 & 0.6335 & 0.5000 & 1187 \\
\texttt{sindy\_multi\_gmz} & 30 & 0.7506 & 0.1227 & -0.0521 & 0.5054 & 0.4720 & 0.5040 & 0.5618 & 0.4570 & 1230 \\
\texttt{sindy\_multi\_gmz} & 45 & 0.9373 & 0.1502 & -0.0480 & 0.3335 & 0.4101 & 0.4458 & 0.5008 & 0.4017 & 1220 \\
\texttt{sindy\_multi\_gmz} & 60 & 0.9320 & 0.1666 & -0.0517 & 0.3203 & 0.3723 & 0.4100 & 0.4646 & 0.3669 & 1201 \\
\addlinespace[1pt]
\texttt{mamba\_single\_idw} & 15 & 0.6642 & 0.1302 & -0.0556 & 0.6434 & 0.5278 & 0.5657 & 0.4786 & 0.6917 & 2198 \\
\texttt{mamba\_single\_idw} & 30 & 0.7216 & 0.1304 & -0.0697 & 0.5900 & 0.5090 & 0.5471 & 0.4783 & 0.6391 & 2032 \\
\texttt{mamba\_single\_idw} & 45 & 0.7717 & 0.1462 & -0.0621 & 0.4544 & 0.4282 & 0.4764 & 0.3845 & 0.6259 & 2476 \\
\texttt{mamba\_single\_idw} & 60 & 0.7796 & 0.1539 & -0.0653 & 0.4054 & 0.3678 & 0.4233 & 0.3279 & 0.5970 & 2769 \\
\addlinespace[1pt]
\texttt{mamba\_multi\_idw} & 15 & 0.6986 & 0.1239 & -0.0505 & 0.6197 & 0.5362 & 0.5731 & 0.4842 & 0.7021 & 2181 \\
\texttt{mamba\_multi\_idw} & 30 & 0.7296 & 0.1316 & -0.0667 & 0.5848 & 0.4770 & 0.5196 & 0.4316 & 0.6528 & 2287 \\
\texttt{mamba\_multi\_idw} & 45 & 0.7566 & 0.1419 & -0.0630 & 0.4933 & 0.4211 & 0.4701 & 0.3778 & 0.6220 & 2504 \\
\texttt{mamba\_multi\_idw} & 60 & 0.7716 & 0.1540 & -0.0540 & 0.4299 & 0.3599 & 0.4168 & 0.3190 & 0.6009 & 2865 \\
\addlinespace[1pt]
\texttt{mamba\_single\_gmz} & 15 & 0.6755 & 0.1151 & -0.0644 & 0.6432 & 0.5331 & 0.5624 & 0.6087 & 0.5227 & 1306 \\
\texttt{mamba\_single\_gmz} & 30 & 0.7093 & 0.1248 & -0.0733 & 0.5807 & 0.4956 & 0.5276 & 0.5664 & 0.4938 & 1326 \\
\texttt{mamba\_single\_gmz} & 45 & 0.7448 & 0.1347 & -0.0661 & 0.4974 & 0.4559 & 0.4935 & 0.4845 & 0.5030 & 1579 \\
\texttt{mamba\_single\_gmz} & 60 & 0.7669 & 0.1416 & -0.0683 & 0.4620 & 0.4010 & 0.4424 & 0.4354 & 0.4497 & 1571 \\
\addlinespace[1pt]
\texttt{mamba\_multi\_gmz} & 15 & 0.6716 & 0.1130 & -0.0683 & 0.6538 & 0.5424 & 0.5713 & 0.6075 & 0.5392 & 1335 \\
\texttt{mamba\_multi\_gmz} & 30 & 0.7152 & 0.1262 & -0.0713 & 0.5869 & 0.4805 & 0.5134 & 0.5489 & 0.4821 & 1328 \\
\texttt{mamba\_multi\_gmz} & 45 & 0.7550 & 0.1393 & -0.0711 & 0.4802 & 0.4361 & 0.4726 & 0.4971 & 0.4504 & 1378 \\
\texttt{mamba\_multi\_gmz} & 60 & 0.7705 & 0.1471 & -0.0686 & 0.4326 & 0.4171 & 0.4549 & 0.4773 & 0.4346 & 1385 \\
\addlinespace[1pt]
\texttt{transformer\_idw\_fullperiod} & 15 & 0.7824 & 0.1864 & -0.0283 & 0.3989 & 0.2173 & 0.3018 & 0.1899 & 0.7350 & 5887 \\
\texttt{transformer\_idw\_fullperiod} & 30 & 0.7942 & 0.2021 & -0.0188 & 0.3916 & 0.1272 & 0.2281 & 0.1341 & 0.7620 & 8642 \\
\texttt{transformer\_idw\_fullperiod} & 45 & 0.8045 & 0.2124 & -0.0087 & 0.3366 & 0.1188 & 0.2212 & 0.1293 & 0.7633 & 8976 \\
\texttt{transformer\_idw\_fullperiod} & 60 & 0.8117 & 0.2189 & -0.0056 & 0.2908 & 0.0939 & 0.2014 & 0.1155 & 0.7870 & 10366 \\
\addlinespace[1pt]
\texttt{transformer\_gmz\_fullperiod} & 15 & 0.7868 & 0.1965 & -0.0169 & 0.3631 & 0.1207 & 0.2214 & 0.1313 & 0.7041 & 8154 \\
\texttt{transformer\_gmz\_fullperiod} & 30 & 0.8020 & 0.2323 & 0.0120 & 0.3171 & 0.0607 & 0.1738 & 0.0981 & 0.7640 & 11850 \\
\texttt{transformer\_gmz\_fullperiod} & 45 & 0.8116 & 0.2225 & -0.0029 & 0.2663 & 0.0932 & 0.1996 & 0.1157 & 0.7285 & 9580 \\
\texttt{transformer\_gmz\_fullperiod} & 60 & 0.8206 & 0.2272 & -0.0090 & 0.2143 & 0.0815 & 0.1899 & 0.1094 & 0.7193 & 9998 \\
\addlinespace[1pt]
\texttt{sindy\_idw\_fullperiod} & 15 & 0.8516 & 0.2021 & -0.0086 & 0.2377 & 0.1998 & 0.2855 & 0.1810 & 0.6746 & 5667 \\
\texttt{sindy\_idw\_fullperiod} & 30 & 1.2977 & 0.2340 & 0.0197 & 0.1204 & 0.1432 & 0.2377 & 0.1465 & 0.6292 & 6532 \\
\texttt{sindy\_idw\_fullperiod} & 45 & 1.2901 & 0.2437 & 0.0261 & 0.1455 & 0.1079 & 0.2085 & 0.1257 & 0.6114 & 7400 \\
\texttt{sindy\_idw\_fullperiod} & 60 & 1.2923 & 0.2626 & 0.0365 & 0.1276 & 0.0817 & 0.1871 & 0.1109 & 0.5989 & 8216 \\
\addlinespace[1pt]
\texttt{sindy\_gmz\_fullperiod} & 15 & 0.8837 & 0.2044 & -0.0037 & 0.2346 & 0.1885 & 0.2745 & 0.1760 & 0.6239 & 5393 \\
\texttt{sindy\_gmz\_fullperiod} & 30 & 20.6629 & 0.5707 & 0.3606 & 0.0051 & 0.1882 & 0.2723 & 0.1786 & 0.5720 & 4870 \\
\texttt{sindy\_gmz\_fullperiod} & 45 & 9.9520 & 0.4395 & 0.2265 & 0.2478 & 0.1881 & 0.2698 & 0.1822 & 0.5201 & 4342 \\
\texttt{sindy\_gmz\_fullperiod} & 60 & 7.9144 & 0.3495 & 0.1221 & 0.0611 & 0.1834 & 0.2635 & 0.1828 & 0.4714 & 3922 \\
\addlinespace[1pt]
\texttt{mamba\_idw\_fullperiod} & 15 & 0.7932 & 0.2000 & -0.0146 & 0.3654 & 0.1397 & 0.2379 & 0.1416 & 0.7436 & 7987 \\
\texttt{mamba\_idw\_fullperiod} & 30 & 0.7795 & 0.1889 & -0.0242 & 0.4147 & 0.1710 & 0.2632 & 0.1608 & 0.7258 & 6867 \\
\texttt{mamba\_idw\_fullperiod} & 45 & 0.8013 & 0.2054 & -0.0140 & 0.3357 & 0.1390 & 0.2371 & 0.1415 & 0.7318 & 7866 \\
\texttt{mamba\_idw\_fullperiod} & 60 & 0.8147 & 0.1990 & -0.0307 & 0.2909 & 0.1538 & 0.2482 & 0.1513 & 0.6903 & 6941 \\
\addlinespace[1pt]
\texttt{mamba\_gmz\_fullperiod} & 15 & 0.7895 & 0.2001 & -0.0037 & 0.3396 & 0.1400 & 0.2368 & 0.1429 & 0.6903 & 7348 \\
\texttt{mamba\_gmz\_fullperiod} & 30 & 0.7945 & 0.2145 & 0.0021 & 0.3261 & 0.1178 & 0.2186 & 0.1300 & 0.6877 & 8048 \\
\texttt{mamba\_gmz\_fullperiod} & 45 & 0.8033 & 0.2138 & -0.0063 & 0.2978 & 0.1191 & 0.2193 & 0.1311 & 0.6700 & 7772 \\
\texttt{mamba\_gmz\_fullperiod} & 60 & 0.8130 & 0.2132 & -0.0137 & 0.2563 & 0.1206 & 0.2199 & 0.1325 & 0.6476 & 7436 \\
\addlinespace[1pt]
\texttt{pysteps} & 15 & 0.6308 & 0.1009 & 0.0042 & 0.7128 & 0.7010 & 0.7226 & 0.6736 & 0.7793 & 1740 \\
\texttt{pysteps} & 30 & 0.7385 & 0.1184 & 0.0050 & 0.6342 & 0.6596 & 0.6836 & 0.6565 & 0.7130 & 1642 \\
\texttt{pysteps} & 45 & 0.7996 & 0.1285 & -0.0096 & 0.5688 & 0.5534 & 0.5806 & 0.6470 & 0.5266 & 1238 \\
\texttt{pysteps} & 60 & 0.8679 & 0.1348 & -0.0567 & 0.3986 & 0.3668 & 0.3937 & 0.6313 & 0.2860 & 689 \\
\end{longtable}
}

\section{Full evaluation metrics versus radar (outer crop)}\label{app:radar_outer}

{\footnotesize
\begin{longtable}{p{4.0cm}cccccccccc}
\caption{Full evaluation metrics versus the SMHI radar field (outer crop $57.55$--$57.85\,^\circ$N $\times$ $11.70$--$12.20\,^\circ$E) on the held-out test window. Same column convention as Table~\ref{tab:app_gauges}. The number of observed wet pixels at the radar sample is $\approx 52985$.}\label{tab:app_radar_outer} \\
\toprule
Model & $h$ & RMSE & MAE & Bias & CC & HSS & F1 & Prec. & Rec. & $n_{\text{wet}}^{\text{pred}}$ \\
\midrule
\endfirsthead
\multicolumn{11}{c}{\tablename\ \thetable{} --- \textit{continued from previous page}} \\
\toprule
Model & $h$ & RMSE & MAE & Bias & CC & HSS & F1 & Prec. & Rec. & $n_{\text{wet}}^{\text{pred}}$ \\
\midrule
\endhead
\midrule
\multicolumn{11}{r}{\textit{continued on next page}} \\
\endfoot
\bottomrule
\endlastfoot
\texttt{naive\_radar} & 15 & 0.6134 & 0.0924 & -0.0005 & 0.7033 & 0.7420 & 0.7621 & 0.7641 & 0.7601 & 52441 \\
\texttt{naive\_radar} & 30 & 0.7412 & 0.1224 & -0.0011 & 0.5671 & 0.6291 & 0.6580 & 0.6614 & 0.6546 & 52441 \\
\texttt{naive\_radar} & 45 & 0.8034 & 0.1381 & -0.0015 & 0.4915 & 0.5611 & 0.5952 & 0.5999 & 0.5906 & 52169 \\
\texttt{naive\_radar} & 60 & 0.8428 & 0.1477 & -0.0020 & 0.4405 & 0.5112 & 0.5491 & 0.5549 & 0.5435 & 51897 \\
\addlinespace[1pt]
\texttt{naive\_gauge} & 15 & 0.6082 & 0.1000 & 0.0053 & 0.6943 & 0.6836 & 0.7078 & 0.7206 & 0.6955 & 50874 \\
\texttt{naive\_gauge} & 30 & 0.7133 & 0.1219 & 0.0047 & 0.5794 & 0.6302 & 0.6586 & 0.6723 & 0.6455 & 50874 \\
\texttt{naive\_gauge} & 45 & 0.7873 & 0.1399 & 0.0040 & 0.4872 & 0.5642 & 0.5975 & 0.6116 & 0.5841 & 50602 \\
\texttt{naive\_gauge} & 60 & 0.8258 & 0.1515 & 0.0031 & 0.4350 & 0.5070 & 0.5446 & 0.5590 & 0.5310 & 50330 \\
\addlinespace[1pt]
\texttt{naive\_cml} & 15 & 0.6372 & 0.1117 & -0.0464 & 0.6061 & 0.5487 & 0.5822 & 0.6122 & 0.5549 & 47773 \\
\texttt{naive\_cml} & 30 & 0.6374 & 0.1127 & -0.0471 & 0.6066 & 0.5252 & 0.5605 & 0.5910 & 0.5329 & 47773 \\
\texttt{naive\_cml} & 45 & 0.6967 & 0.1275 & -0.0471 & 0.4925 & 0.4607 & 0.5007 & 0.5280 & 0.4761 & 47773 \\
\texttt{naive\_cml} & 60 & 0.7327 & 0.1386 & -0.0471 & 0.4183 & 0.4021 & 0.4465 & 0.4709 & 0.4245 & 47773 \\
\addlinespace[1pt]
\texttt{model1\_physical} & 15 & 0.7362 & 0.1757 & -0.0107 & 0.3849 & 0.2744 & 0.3567 & 0.2450 & 0.6554 & 141003 \\
\texttt{model1\_physical} & 30 & 0.7368 & 0.1762 & -0.0112 & 0.3847 & 0.2759 & 0.3584 & 0.2464 & 0.6570 & 141262 \\
\texttt{model1\_physical} & 45 & 0.7370 & 0.1763 & -0.0112 & 0.3847 & 0.2759 & 0.3584 & 0.2464 & 0.6570 & 141262 \\
\texttt{model1\_physical} & 60 & 0.7371 & 0.1763 & -0.0113 & 0.3847 & 0.2759 & 0.3584 & 0.2464 & 0.6570 & 141262 \\
\addlinespace[1pt]
\texttt{transformer\_single\_idw} & 15 & 0.6601 & 0.1296 & -0.0540 & 0.5671 & 0.5426 & 0.5784 & 0.5764 & 0.5804 & 53072 \\
\texttt{transformer\_single\_idw} & 30 & 0.6607 & 0.1249 & -0.0692 & 0.6210 & 0.4885 & 0.5321 & 0.4861 & 0.5878 & 63752 \\
\texttt{transformer\_single\_idw} & 45 & 0.7071 & 0.1700 & -0.1031 & 0.5299 & 0.4203 & 0.4736 & 0.4018 & 0.5765 & 75636 \\
\texttt{transformer\_single\_idw} & 60 & 0.7436 & 0.1565 & -0.0496 & 0.4000 & 0.3422 & 0.4074 & 0.3203 & 0.5596 & 92075 \\
\addlinespace[1pt]
\texttt{transformer\_multi\_idw} & 15 & 0.6520 & 0.1212 & -0.0537 & 0.6117 & 0.5171 & 0.5579 & 0.5075 & 0.6193 & 63326 \\
\texttt{transformer\_multi\_idw} & 30 & 0.6549 & 0.1252 & -0.0556 & 0.6193 & 0.4686 & 0.5159 & 0.4475 & 0.6088 & 70970 \\
\texttt{transformer\_multi\_idw} & 45 & 0.6786 & 0.1358 & -0.0590 & 0.5650 & 0.4129 & 0.4667 & 0.3947 & 0.5709 & 75840 \\
\texttt{transformer\_multi\_idw} & 60 & 0.7030 & 0.1485 & -0.0708 & 0.5042 & 0.3689 & 0.4268 & 0.3628 & 0.5184 & 75326 \\
\addlinespace[1pt]
\texttt{transformer\_single\_gmz} & 15 & 0.7772 & 0.1294 & -0.1150 & 0.2640 & 0.1371 & 0.1559 & 0.5557 & 0.0907 & 8601 \\
\texttt{transformer\_single\_gmz} & 30 & 0.7812 & 0.1299 & -0.1166 & 0.2739 & 0.1549 & 0.1757 & 0.5568 & 0.1043 & 9875 \\
\texttt{transformer\_single\_gmz} & 45 & 0.7860 & 0.1336 & -0.1094 & 0.2266 & 0.1415 & 0.1644 & 0.4809 & 0.0992 & 10869 \\
\texttt{transformer\_single\_gmz} & 60 & 0.7906 & 0.1336 & -0.1141 & 0.2034 & 0.1277 & 0.1496 & 0.4639 & 0.0892 & 10136 \\
\addlinespace[1pt]
\texttt{transformer\_multi\_gmz} & 15 & 0.7848 & 0.1300 & -0.1133 & 0.2362 & 0.1779 & 0.2001 & 0.5676 & 0.1215 & 11107 \\
\texttt{transformer\_multi\_gmz} & 30 & 0.7844 & 0.1296 & -0.1130 & 0.2504 & 0.1646 & 0.1870 & 0.5396 & 0.1131 & 10932 \\
\texttt{transformer\_multi\_gmz} & 45 & 0.7873 & 0.1318 & -0.1119 & 0.2319 & 0.1486 & 0.1723 & 0.4807 & 0.1049 & 11449 \\
\texttt{transformer\_multi\_gmz} & 60 & 0.7905 & 0.1334 & -0.1126 & 0.2042 & 0.1353 & 0.1596 & 0.4442 & 0.0972 & 11540 \\
\addlinespace[1pt]
\texttt{sindy\_single\_idw} & 15 & 0.7084 & 0.1256 & -0.0491 & 0.4907 & 0.5072 & 0.5436 & 0.5721 & 0.5178 & 47711 \\
\texttt{sindy\_single\_idw} & 30 & 0.8613 & 0.1514 & -0.0502 & 0.3591 & 0.4371 & 0.4790 & 0.5016 & 0.4583 & 48161 \\
\texttt{sindy\_single\_idw} & 45 & 7.0837 & 0.3996 & 0.0784 & 0.0114 & 0.3495 & 0.4003 & 0.3997 & 0.4008 & 52856 \\
\texttt{sindy\_single\_idw} & 60 & 1.9125 & 0.2416 & -0.0610 & 0.0837 & 0.2918 & 0.3507 & 0.3287 & 0.3758 & 60274 \\
\addlinespace[1pt]
\texttt{sindy\_multi\_idw} & 15 & 0.6430 & 0.1162 & -0.0476 & 0.5942 & 0.5190 & 0.5547 & 0.5736 & 0.5370 & 48585 \\
\texttt{sindy\_multi\_idw} & 30 & 0.6635 & 0.1270 & -0.0510 & 0.5584 & 0.4583 & 0.4991 & 0.5113 & 0.4874 & 49729 \\
\texttt{sindy\_multi\_idw} & 45 & 0.7533 & 0.1529 & -0.0512 & 0.4197 & 0.3677 & 0.4176 & 0.4103 & 0.4253 & 54356 \\
\texttt{sindy\_multi\_idw} & 60 & 1.9125 & 0.2416 & -0.0610 & 0.0837 & 0.2918 & 0.3507 & 0.3287 & 0.3758 & 60274 \\
\addlinespace[1pt]
\texttt{sindy\_single\_gmz} & 15 & 0.7735 & 0.1238 & -0.1109 & 0.2769 & 0.1380 & 0.1554 & 0.5984 & 0.0893 & 7868 \\
\texttt{sindy\_single\_gmz} & 30 & 0.7761 & 0.1244 & -0.1102 & 0.2743 & 0.1320 & 0.1499 & 0.5684 & 0.0863 & 8005 \\
\texttt{sindy\_single\_gmz} & 45 & 0.8073 & 0.1291 & -0.1096 & 0.2083 & 0.1127 & 0.1309 & 0.4974 & 0.0754 & 7989 \\
\texttt{sindy\_single\_gmz} & 60 & 0.8192 & 0.1331 & -0.1107 & 0.1833 & 0.0986 & 0.1168 & 0.4511 & 0.0671 & 7838 \\
\addlinespace[1pt]
\texttt{sindy\_multi\_gmz} & 15 & 0.7727 & 0.1226 & -0.1096 & 0.2773 & 0.1402 & 0.1576 & 0.5984 & 0.0907 & 7868 \\
\texttt{sindy\_multi\_gmz} & 30 & 0.7757 & 0.1236 & -0.1093 & 0.2746 & 0.1334 & 0.1512 & 0.5684 & 0.0872 & 8005 \\
\texttt{sindy\_multi\_gmz} & 45 & 0.8071 & 0.1286 & -0.1092 & 0.2084 & 0.1133 & 0.1315 & 0.4974 & 0.0758 & 7989 \\
\texttt{sindy\_multi\_gmz} & 60 & 0.8192 & 0.1331 & -0.1107 & 0.1833 & 0.0986 & 0.1168 & 0.4511 & 0.0671 & 7838 \\
\addlinespace[1pt]
\texttt{mamba\_single\_idw} & 15 & 0.6605 & 0.1363 & -0.0537 & 0.5694 & 0.5030 & 0.5452 & 0.5005 & 0.5987 & 63054 \\
\texttt{mamba\_single\_idw} & 30 & 0.6798 & 0.1264 & -0.0664 & 0.5847 & 0.4873 & 0.5288 & 0.5087 & 0.5506 & 57063 \\
\texttt{mamba\_single\_idw} & 45 & 0.7196 & 0.1379 & -0.0589 & 0.4960 & 0.4223 & 0.4739 & 0.4138 & 0.5545 & 70632 \\
\texttt{mamba\_single\_idw} & 60 & 0.7343 & 0.1459 & -0.0631 & 0.4262 & 0.3744 & 0.4332 & 0.3594 & 0.5452 & 79954 \\
\addlinespace[1pt]
\texttt{mamba\_multi\_idw} & 15 & 0.6719 & 0.1275 & -0.0473 & 0.5781 & 0.5106 & 0.5514 & 0.5087 & 0.6018 & 61401 \\
\texttt{mamba\_multi\_idw} & 30 & 0.6854 & 0.1264 & -0.0595 & 0.5866 & 0.4756 & 0.5196 & 0.4777 & 0.5696 & 62200 \\
\texttt{mamba\_multi\_idw} & 45 & 0.7039 & 0.1345 & -0.0571 & 0.5344 & 0.4184 & 0.4704 & 0.4081 & 0.5552 & 71333 \\
\texttt{mamba\_multi\_idw} & 60 & 0.7205 & 0.1456 & -0.0504 & 0.4636 & 0.3656 & 0.4263 & 0.3477 & 0.5507 & 83482 \\
\addlinespace[1pt]
\texttt{mamba\_single\_gmz} & 15 & 0.7765 & 0.1265 & -0.1116 & 0.2707 & 0.1562 & 0.1777 & 0.5431 & 0.1062 & 10311 \\
\texttt{mamba\_single\_gmz} & 30 & 0.7753 & 0.1271 & -0.1129 & 0.2839 & 0.1447 & 0.1659 & 0.5224 & 0.0986 & 9951 \\
\texttt{mamba\_single\_gmz} & 45 & 0.7809 & 0.1285 & -0.1127 & 0.2601 & 0.1400 & 0.1636 & 0.4656 & 0.0992 & 11233 \\
\texttt{mamba\_single\_gmz} & 60 & 0.7877 & 0.1308 & -0.1128 & 0.2290 & 0.1196 & 0.1443 & 0.4012 & 0.0880 & 11564 \\
\addlinespace[1pt]
\texttt{mamba\_multi\_gmz} & 15 & 0.7756 & 0.1251 & -0.1112 & 0.2734 & 0.1553 & 0.1764 & 0.5426 & 0.1053 & 10073 \\
\texttt{mamba\_multi\_gmz} & 30 & 0.7751 & 0.1261 & -0.1120 & 0.2900 & 0.1437 & 0.1652 & 0.5063 & 0.0987 & 10171 \\
\texttt{mamba\_multi\_gmz} & 45 & 0.7826 & 0.1289 & -0.1128 & 0.2507 & 0.1271 & 0.1491 & 0.4583 & 0.0891 & 10189 \\
\texttt{mamba\_multi\_gmz} & 60 & 0.7866 & 0.1312 & -0.1126 & 0.2265 & 0.1194 & 0.1422 & 0.4277 & 0.0853 & 10514 \\
\addlinespace[1pt]
\texttt{transformer\_idw\_fullperiod} & 15 & 0.7360 & 0.1794 & -0.0210 & 0.4227 & 0.2043 & 0.3001 & 0.1942 & 0.6601 & 179166 \\
\texttt{transformer\_idw\_fullperiod} & 30 & 0.7544 & 0.1976 & -0.0099 & 0.3736 & 0.1128 & 0.2282 & 0.1367 & 0.6917 & 266774 \\
\texttt{transformer\_idw\_fullperiod} & 45 & 0.7626 & 0.2151 & 0.0068 & 0.3389 & 0.0889 & 0.2095 & 0.1234 & 0.6950 & 296969 \\
\texttt{transformer\_idw\_fullperiod} & 60 & 0.7715 & 0.2192 & 0.0079 & 0.2737 & 0.0685 & 0.1938 & 0.1123 & 0.7040 & 330358 \\
\addlinespace[1pt]
\texttt{transformer\_gmz\_fullperiod} & 15 & 0.7994 & 0.1424 & -0.1055 & 0.1057 & 0.0893 & 0.1463 & 0.1821 & 0.1222 & 35369 \\
\texttt{transformer\_gmz\_fullperiod} & 30 & 0.8021 & 0.1440 & -0.1059 & 0.0753 & 0.0718 & 0.1355 & 0.1539 & 0.1210 & 41448 \\
\texttt{transformer\_gmz\_fullperiod} & 45 & 0.8028 & 0.1432 & -0.1069 & 0.0679 & 0.0667 & 0.1275 & 0.1526 & 0.1095 & 37823 \\
\texttt{transformer\_gmz\_fullperiod} & 60 & 0.8042 & 0.1434 & -0.1082 & 0.0519 & 0.0674 & 0.1276 & 0.1542 & 0.1087 & 37164 \\
\addlinespace[1pt]
\texttt{sindy\_idw\_fullperiod} & 15 & 0.7764 & 0.1926 & 5.91e-05 & 0.3039 & 0.1909 & 0.2875 & 0.1875 & 0.6166 & 173364 \\
\texttt{sindy\_idw\_fullperiod} & 30 & 1.0544 & 0.2249 & 0.0280 & 0.1773 & 0.1161 & 0.2271 & 0.1409 & 0.5854 & 218989 \\
\texttt{sindy\_idw\_fullperiod} & 45 & 1.1398 & 0.2374 & 0.0373 & 0.1907 & 0.0697 & 0.1903 & 0.1145 & 0.5645 & 259979 \\
\texttt{sindy\_idw\_fullperiod} & 60 & 1.1134 & 0.2576 & 0.0498 & 0.1822 & 0.0475 & 0.1730 & 0.1024 & 0.5569 & 286688 \\
\addlinespace[1pt]
\texttt{sindy\_gmz\_fullperiod} & 15 & 0.8087 & 0.1391 & -0.1049 & 0.0866 & 0.0942 & 0.1390 & 0.2190 & 0.1018 & 24508 \\
\texttt{sindy\_gmz\_fullperiod} & 30 & 6.7084 & 0.1893 & -0.0580 & 0.0018 & 0.0840 & 0.1266 & 0.2122 & 0.0902 & 22402 \\
\texttt{sindy\_gmz\_fullperiod} & 45 & 3.2851 & 0.1702 & -0.0753 & 0.0974 & 0.0761 & 0.1166 & 0.2074 & 0.0811 & 20626 \\
\texttt{sindy\_gmz\_fullperiod} & 60 & 2.6569 & 0.1588 & -0.0887 & 0.0289 & 0.0676 & 0.1066 & 0.1987 & 0.0728 & 19322 \\
\addlinespace[1pt]
\texttt{mamba\_idw\_fullperiod} & 15 & 0.7466 & 0.1944 & -0.0068 & 0.3874 & 0.1312 & 0.2424 & 0.1475 & 0.6810 & 243436 \\
\texttt{mamba\_idw\_fullperiod} & 30 & 0.7403 & 0.1865 & -0.0132 & 0.4082 & 0.1738 & 0.2758 & 0.1740 & 0.6641 & 201133 \\
\texttt{mamba\_idw\_fullperiod} & 45 & 0.7639 & 0.2008 & -0.0087 & 0.3119 & 0.1236 & 0.2361 & 0.1434 & 0.6676 & 245445 \\
\texttt{mamba\_idw\_fullperiod} & 60 & 0.7746 & 0.1972 & -0.0200 & 0.2660 & 0.1312 & 0.2410 & 0.1487 & 0.6363 & 225585 \\
\addlinespace[1pt]
\texttt{mamba\_gmz\_fullperiod} & 15 & 0.7975 & 0.1407 & -0.1047 & 0.1249 & 0.0926 & 0.1461 & 0.1929 & 0.1176 & 32134 \\
\texttt{mamba\_gmz\_fullperiod} & 30 & 0.7993 & 0.1427 & -0.1035 & 0.1059 & 0.0750 & 0.1354 & 0.1615 & 0.1166 & 38044 \\
\texttt{mamba\_gmz\_fullperiod} & 45 & 0.8010 & 0.1430 & -0.1042 & 0.0893 & 0.0741 & 0.1326 & 0.1637 & 0.1115 & 35908 \\
\texttt{mamba\_gmz\_fullperiod} & 60 & 0.8030 & 0.1417 & -0.1074 & 0.0705 & 0.0691 & 0.1255 & 0.1617 & 0.1025 & 33435 \\
\addlinespace[1pt]
\texttt{pysteps} & 15 & 0.5849 & 0.0881 & 0.0051 & 0.7406 & 0.7760 & 0.7933 & 0.7912 & 0.7953 & 52165 \\
\texttt{pysteps} & 30 & 0.6270 & 0.0973 & 0.0050 & 0.7230 & 0.7239 & 0.7441 & 0.7884 & 0.7044 & 46614 \\
\texttt{pysteps} & 45 & 0.7200 & 0.1114 & -0.0184 & 0.6099 & 0.5903 & 0.6154 & 0.7779 & 0.5091 & 34322 \\
\texttt{pysteps} & 60 & 0.7922 & 0.1232 & -0.0593 & 0.4208 & 0.3913 & 0.4183 & 0.7355 & 0.2923 & 20947 \\
\end{longtable}
}

\section{Full evaluation metrics versus radar (inner crop)}\label{app:radar_inner}

{\footnotesize
\begin{longtable}{p{4.0cm}cccccccccc}
\caption{Full evaluation metrics versus the SMHI radar field, inner crop ($57.6$--$57.8\,^\circ$N $\times$ $11.75$--$12.10\,^\circ$E) on the held-out test window. Same column convention as Table~\ref{tab:app_gauges}. The number of observed wet pixels at the radar sample is $\approx 23404$.}\label{tab:app_radar_inner} \\
\toprule
Model & $h$ & RMSE & MAE & Bias & CC & HSS & F1 & Prec. & Rec. & $n_{\text{wet}}^{\text{pred}}$ \\
\midrule
\endfirsthead
\multicolumn{11}{c}{\tablename\ \thetable{} --- \textit{continued from previous page}} \\
\toprule
Model & $h$ & RMSE & MAE & Bias & CC & HSS & F1 & Prec. & Rec. & $n_{\text{wet}}^{\text{pred}}$ \\
\midrule
\endhead
\midrule
\multicolumn{11}{r}{\textit{continued on next page}} \\
\endfoot
\bottomrule
\endlastfoot
\texttt{naive\_radar} & 15 & 0.6140 & 0.0931 & -0.0005 & 0.7151 & 0.7495 & 0.7690 & 0.7710 & 0.7670 & 23164 \\
\texttt{naive\_radar} & 30 & 0.7475 & 0.1242 & -0.0011 & 0.5780 & 0.6344 & 0.6629 & 0.6663 & 0.6595 & 23164 \\
\texttt{naive\_radar} & 45 & 0.8090 & 0.1404 & -0.0014 & 0.5059 & 0.5694 & 0.6029 & 0.6076 & 0.5983 & 23044 \\
\texttt{naive\_radar} & 60 & 0.8460 & 0.1502 & -0.0019 & 0.4596 & 0.5213 & 0.5585 & 0.5643 & 0.5527 & 22924 \\
\addlinespace[1pt]
\texttt{naive\_gauge} & 15 & 0.5692 & 0.0926 & 0.0033 & 0.7465 & 0.7070 & 0.7292 & 0.7507 & 0.7088 & 21984 \\
\texttt{naive\_gauge} & 30 & 0.7174 & 0.1225 & 0.0027 & 0.5972 & 0.6427 & 0.6699 & 0.6915 & 0.6495 & 21984 \\
\texttt{naive\_gauge} & 45 & 0.8077 & 0.1442 & 0.0021 & 0.4889 & 0.5706 & 0.6032 & 0.6244 & 0.5833 & 21864 \\
\texttt{naive\_gauge} & 60 & 0.8474 & 0.1565 & 0.0012 & 0.4368 & 0.5106 & 0.5476 & 0.5685 & 0.5282 & 21744 \\
\addlinespace[1pt]
\texttt{naive\_cml} & 15 & 0.6221 & 0.1069 & -0.0506 & 0.6575 & 0.5732 & 0.6039 & 0.6534 & 0.5614 & 20005 \\
\texttt{naive\_cml} & 30 & 0.6198 & 0.1083 & -0.0512 & 0.6624 & 0.5441 & 0.5771 & 0.6261 & 0.5352 & 20005 \\
\texttt{naive\_cml} & 45 & 0.7020 & 0.1286 & -0.0512 & 0.5115 & 0.4657 & 0.5043 & 0.5472 & 0.4677 & 20005 \\
\texttt{naive\_cml} & 60 & 0.7442 & 0.1414 & -0.0513 & 0.4269 & 0.4035 & 0.4466 & 0.4845 & 0.4142 & 20005 \\
\addlinespace[1pt]
\texttt{model1\_physical} & 15 & 0.7457 & 0.1760 & -0.0144 & 0.4021 & 0.2718 & 0.3544 & 0.2435 & 0.6513 & 62287 \\
\texttt{model1\_physical} & 30 & 0.7462 & 0.1766 & -0.0150 & 0.4020 & 0.2732 & 0.3560 & 0.2448 & 0.6526 & 62395 \\
\texttt{model1\_physical} & 45 & 0.7464 & 0.1766 & -0.0150 & 0.4020 & 0.2731 & 0.3560 & 0.2448 & 0.6526 & 62395 \\
\texttt{model1\_physical} & 60 & 0.7465 & 0.1767 & -0.0150 & 0.4020 & 0.2731 & 0.3560 & 0.2448 & 0.6526 & 62395 \\
\addlinespace[1pt]
\texttt{transformer\_single\_idw} & 15 & 0.6530 & 0.1270 & -0.0576 & 0.6118 & 0.5682 & 0.6021 & 0.5987 & 0.6055 & 23546 \\
\texttt{transformer\_single\_idw} & 30 & 0.6600 & 0.1245 & -0.0731 & 0.6699 & 0.4983 & 0.5412 & 0.4945 & 0.5975 & 28131 \\
\texttt{transformer\_single\_idw} & 45 & 0.7166 & 0.1724 & -0.1066 & 0.5549 & 0.4234 & 0.4765 & 0.4041 & 0.5803 & 33431 \\
\texttt{transformer\_single\_idw} & 60 & 0.7581 & 0.1592 & -0.0538 & 0.4133 & 0.3479 & 0.4123 & 0.3263 & 0.5597 & 39944 \\
\addlinespace[1pt]
\texttt{transformer\_multi\_idw} & 15 & 0.6523 & 0.1213 & -0.0575 & 0.6508 & 0.5377 & 0.5768 & 0.5245 & 0.6406 & 27995 \\
\texttt{transformer\_multi\_idw} & 30 & 0.6545 & 0.1251 & -0.0593 & 0.6640 & 0.4773 & 0.5239 & 0.4548 & 0.6176 & 31294 \\
\texttt{transformer\_multi\_idw} & 45 & 0.6861 & 0.1381 & -0.0647 & 0.5915 & 0.4146 & 0.4682 & 0.3973 & 0.5698 & 33218 \\
\texttt{transformer\_multi\_idw} & 60 & 0.7142 & 0.1515 & -0.0756 & 0.5214 & 0.3715 & 0.4290 & 0.3663 & 0.5176 & 32898 \\
\addlinespace[1pt]
\texttt{transformer\_single\_gmz} & 15 & 0.7572 & 0.1333 & -0.1041 & 0.3903 & 0.2670 & 0.2977 & 0.5545 & 0.2034 & 8543 \\
\texttt{transformer\_single\_gmz} & 30 & 0.7664 & 0.1324 & -0.1074 & 0.4055 & 0.2974 & 0.3298 & 0.5587 & 0.2339 & 9750 \\
\texttt{transformer\_single\_gmz} & 45 & 0.7772 & 0.1407 & -0.0915 & 0.3385 & 0.2692 & 0.3054 & 0.4810 & 0.2238 & 10832 \\
\texttt{transformer\_single\_gmz} & 60 & 0.7872 & 0.1395 & -0.1012 & 0.3038 & 0.2441 & 0.2798 & 0.4627 & 0.2006 & 10093 \\
\addlinespace[1pt]
\texttt{transformer\_multi\_gmz} & 15 & 0.7760 & 0.1347 & -0.1017 & 0.3503 & 0.3337 & 0.3669 & 0.5660 & 0.2714 & 10991 \\
\texttt{transformer\_multi\_gmz} & 30 & 0.7740 & 0.1325 & -0.0998 & 0.3741 & 0.3104 & 0.3445 & 0.5382 & 0.2533 & 10847 \\
\texttt{transformer\_multi\_gmz} & 45 & 0.7804 & 0.1373 & -0.0967 & 0.3463 & 0.2772 & 0.3139 & 0.4793 & 0.2334 & 11281 \\
\texttt{transformer\_multi\_gmz} & 60 & 0.7874 & 0.1403 & -0.0977 & 0.3025 & 0.2529 & 0.2913 & 0.4428 & 0.2171 & 11413 \\
\addlinespace[1pt]
\texttt{sindy\_single\_idw} & 15 & 0.7027 & 0.1238 & -0.0544 & 0.5276 & 0.5240 & 0.5584 & 0.6026 & 0.5202 & 20099 \\
\texttt{sindy\_single\_idw} & 30 & 0.8606 & 0.1495 & -0.0588 & 0.3759 & 0.4511 & 0.4910 & 0.5274 & 0.4593 & 20278 \\
\texttt{sindy\_single\_idw} & 45 & 6.9081 & 0.3970 & 0.0641 & 0.0039 & 0.3481 & 0.3981 & 0.4050 & 0.3916 & 22513 \\
\texttt{sindy\_single\_idw} & 60 & 1.9549 & 0.2474 & -0.0679 & 0.0901 & 0.2903 & 0.3487 & 0.3308 & 0.3687 & 25949 \\
\addlinespace[1pt]
\texttt{sindy\_multi\_idw} & 15 & 0.6280 & 0.1127 & -0.0520 & 0.6459 & 0.5390 & 0.5725 & 0.6061 & 0.5424 & 20518 \\
\texttt{sindy\_multi\_idw} & 30 & 0.6552 & 0.1247 & -0.0559 & 0.5968 & 0.4758 & 0.5144 & 0.5397 & 0.4914 & 20980 \\
\texttt{sindy\_multi\_idw} & 45 & 0.7707 & 0.1562 & -0.0565 & 0.4242 & 0.3680 & 0.4171 & 0.4172 & 0.4170 & 23155 \\
\texttt{sindy\_multi\_idw} & 60 & 1.9549 & 0.2474 & -0.0679 & 0.0901 & 0.2903 & 0.3487 & 0.3308 & 0.3687 & 25949 \\
\addlinespace[1pt]
\texttt{sindy\_single\_gmz} & 15 & 0.7486 & 0.1238 & -0.0944 & 0.4089 & 0.2736 & 0.3023 & 0.5984 & 0.2022 & 7868 \\
\texttt{sindy\_single\_gmz} & 30 & 0.7546 & 0.1250 & -0.0927 & 0.4049 & 0.2613 & 0.2908 & 0.5684 & 0.1954 & 8005 \\
\texttt{sindy\_single\_gmz} & 45 & 0.8253 & 0.1356 & -0.0915 & 0.3072 & 0.2232 & 0.2541 & 0.4974 & 0.1707 & 7989 \\
\texttt{sindy\_single\_gmz} & 60 & 0.8512 & 0.1445 & -0.0938 & 0.2702 & 0.1956 & 0.2272 & 0.4511 & 0.1519 & 7838 \\
\addlinespace[1pt]
\texttt{sindy\_multi\_gmz} & 15 & 0.7480 & 0.1226 & -0.0932 & 0.4094 & 0.2774 & 0.3058 & 0.5984 & 0.2054 & 7868 \\
\texttt{sindy\_multi\_gmz} & 30 & 0.7543 & 0.1243 & -0.0919 & 0.4052 & 0.2637 & 0.2931 & 0.5684 & 0.1974 & 8005 \\
\texttt{sindy\_multi\_gmz} & 45 & 0.8251 & 0.1351 & -0.0910 & 0.3074 & 0.2242 & 0.2551 & 0.4974 & 0.1716 & 7989 \\
\texttt{sindy\_multi\_gmz} & 60 & 0.8512 & 0.1445 & -0.0938 & 0.2702 & 0.1956 & 0.2272 & 0.4511 & 0.1519 & 7838 \\
\addlinespace[1pt]
\texttt{mamba\_single\_idw} & 15 & 0.6568 & 0.1356 & -0.0581 & 0.6090 & 0.5244 & 0.5646 & 0.5210 & 0.6162 & 27540 \\
\texttt{mamba\_single\_idw} & 30 & 0.6832 & 0.1254 & -0.0700 & 0.6274 & 0.5076 & 0.5474 & 0.5291 & 0.5669 & 24948 \\
\texttt{mamba\_single\_idw} & 45 & 0.7328 & 0.1391 & -0.0647 & 0.5184 & 0.4272 & 0.4779 & 0.4215 & 0.5516 & 30469 \\
\texttt{mamba\_single\_idw} & 60 & 0.7487 & 0.1486 & -0.0661 & 0.4374 & 0.3769 & 0.4354 & 0.3623 & 0.5455 & 35063 \\
\addlinespace[1pt]
\texttt{mamba\_multi\_idw} & 15 & 0.6776 & 0.1255 & -0.0545 & 0.6141 & 0.5313 & 0.5697 & 0.5341 & 0.6104 & 26199 \\
\texttt{mamba\_multi\_idw} & 30 & 0.6918 & 0.1259 & -0.0636 & 0.6268 & 0.4881 & 0.5314 & 0.4860 & 0.5861 & 27786 \\
\texttt{mamba\_multi\_idw} & 45 & 0.7152 & 0.1364 & -0.0609 & 0.5554 & 0.4197 & 0.4716 & 0.4093 & 0.5563 & 31478 \\
\texttt{mamba\_multi\_idw} & 60 & 0.7340 & 0.1485 & -0.0541 & 0.4755 & 0.3661 & 0.4267 & 0.3487 & 0.5495 & 36687 \\
\addlinespace[1pt]
\texttt{mamba\_single\_gmz} & 15 & 0.7558 & 0.1281 & -0.0960 & 0.3999 & 0.2848 & 0.3183 & 0.5309 & 0.2272 & 9966 \\
\texttt{mamba\_single\_gmz} & 30 & 0.7528 & 0.1295 & -0.0990 & 0.4194 & 0.2661 & 0.2997 & 0.5103 & 0.2121 & 9678 \\
\texttt{mamba\_single\_gmz} & 45 & 0.7659 & 0.1325 & -0.0985 & 0.3839 & 0.2499 & 0.2872 & 0.4515 & 0.2105 & 10856 \\
\texttt{mamba\_single\_gmz} & 60 & 0.7803 & 0.1368 & -0.0983 & 0.3443 & 0.2187 & 0.2586 & 0.3955 & 0.1921 & 11308 \\
\addlinespace[1pt]
\texttt{mamba\_multi\_gmz} & 15 & 0.7548 & 0.1268 & -0.0970 & 0.4040 & 0.2873 & 0.3202 & 0.5341 & 0.2286 & 9813 \\
\texttt{mamba\_multi\_gmz} & 30 & 0.7529 & 0.1283 & -0.0979 & 0.4293 & 0.2648 & 0.2990 & 0.4968 & 0.2139 & 9919 \\
\texttt{mamba\_multi\_gmz} & 45 & 0.7695 & 0.1340 & -0.0992 & 0.3728 & 0.2342 & 0.2700 & 0.4489 & 0.1930 & 9960 \\
\texttt{mamba\_multi\_gmz} & 60 & 0.7786 & 0.1386 & -0.0982 & 0.3345 & 0.2172 & 0.2545 & 0.4167 & 0.1832 & 10237 \\
\addlinespace[1pt]
\texttt{transformer\_idw\_fullperiod} & 15 & 0.7515 & 0.1828 & -0.0237 & 0.4289 & 0.2047 & 0.3004 & 0.1947 & 0.6573 & 78620 \\
\texttt{transformer\_idw\_fullperiod} & 30 & 0.7703 & 0.1969 & -0.0189 & 0.3810 & 0.1215 & 0.2349 & 0.1420 & 0.6808 & 111669 \\
\texttt{transformer\_idw\_fullperiod} & 45 & 0.7787 & 0.2151 & -0.0006 & 0.3421 & 0.0954 & 0.2144 & 0.1272 & 0.6833 & 125106 \\
\texttt{transformer\_idw\_fullperiod} & 60 & 0.7878 & 0.2192 & 0.0004 & 0.2761 & 0.0758 & 0.1994 & 0.1164 & 0.6928 & 138531 \\
\addlinespace[1pt]
\texttt{transformer\_gmz\_fullperiod} & 15 & 0.8081 & 0.1609 & -0.0824 & 0.1583 & 0.1391 & 0.2200 & 0.1826 & 0.2766 & 35274 \\
\texttt{transformer\_gmz\_fullperiod} & 30 & 0.8137 & 0.1646 & -0.0829 & 0.1135 & 0.1082 & 0.1972 & 0.1541 & 0.2740 & 41396 \\
\texttt{transformer\_gmz\_fullperiod} & 45 & 0.8151 & 0.1636 & -0.0851 & 0.1023 & 0.1019 & 0.1887 & 0.1524 & 0.2476 & 37816 \\
\texttt{transformer\_gmz\_fullperiod} & 60 & 0.8179 & 0.1636 & -0.0882 & 0.0786 & 0.1037 & 0.1897 & 0.1543 & 0.2462 & 37142 \\
\addlinespace[1pt]
\texttt{sindy\_idw\_fullperiod} & 15 & 0.8003 & 0.1972 & -0.0036 & 0.2907 & 0.1781 & 0.2773 & 0.1792 & 0.6129 & 79658 \\
\texttt{sindy\_idw\_fullperiod} & 30 & 1.0907 & 0.2293 & 0.0224 & 0.1616 & 0.1085 & 0.2210 & 0.1367 & 0.5773 & 98367 \\
\texttt{sindy\_idw\_fullperiod} & 45 & 1.1464 & 0.2430 & 0.0326 & 0.1850 & 0.0682 & 0.1890 & 0.1138 & 0.5588 & 114370 \\
\texttt{sindy\_idw\_fullperiod} & 60 & 1.1296 & 0.2642 & 0.0451 & 0.1755 & 0.0476 & 0.1730 & 0.1026 & 0.5511 & 125086 \\
\addlinespace[1pt]
\texttt{sindy\_gmz\_fullperiod} & 15 & 0.8287 & 0.1582 & -0.0808 & 0.1280 & 0.1572 & 0.2246 & 0.2190 & 0.2305 & 24508 \\
\texttt{sindy\_gmz\_fullperiod} & 30 & 10.0603 & 0.2721 & 0.0256 & 0.0025 & 0.1425 & 0.2081 & 0.2122 & 0.2042 & 22402 \\
\texttt{sindy\_gmz\_fullperiod} & 45 & 4.8646 & 0.2288 & -0.0135 & 0.1435 & 0.1311 & 0.1948 & 0.2074 & 0.1837 & 20626 \\
\texttt{sindy\_gmz\_fullperiod} & 60 & 3.8993 & 0.2029 & -0.0439 & 0.0425 & 0.1177 & 0.1802 & 0.1987 & 0.1649 & 19322 \\
\addlinespace[1pt]
\texttt{mamba\_idw\_fullperiod} & 15 & 0.7623 & 0.1941 & -0.0145 & 0.3957 & 0.1407 & 0.2498 & 0.1535 & 0.6703 & 101677 \\
\texttt{mamba\_idw\_fullperiod} & 30 & 0.7562 & 0.1894 & -0.0181 & 0.4126 & 0.1708 & 0.2732 & 0.1725 & 0.6555 & 88451 \\
\texttt{mamba\_idw\_fullperiod} & 45 & 0.7807 & 0.2029 & -0.0140 & 0.3116 & 0.1265 & 0.2382 & 0.1454 & 0.6583 & 105425 \\
\texttt{mamba\_idw\_fullperiod} & 60 & 0.7911 & 0.1983 & -0.0259 & 0.2678 & 0.1373 & 0.2454 & 0.1529 & 0.6225 & 94829 \\
\addlinespace[1pt]
\texttt{mamba\_gmz\_fullperiod} & 15 & 0.8037 & 0.1585 & -0.0803 & 0.1872 & 0.1434 & 0.2207 & 0.1906 & 0.2620 & 32002 \\
\texttt{mamba\_gmz\_fullperiod} & 30 & 0.8075 & 0.1638 & -0.0775 & 0.1595 & 0.1129 & 0.1987 & 0.1603 & 0.2613 & 37956 \\
\texttt{mamba\_gmz\_fullperiod} & 45 & 0.8110 & 0.1647 & -0.0791 & 0.1351 & 0.1135 & 0.1974 & 0.1628 & 0.2506 & 35845 \\
\texttt{mamba\_gmz\_fullperiod} & 60 & 0.8154 & 0.1619 & -0.0864 & 0.1058 & 0.1076 & 0.1897 & 0.1610 & 0.2308 & 33391 \\
\addlinespace[1pt]
\texttt{pysteps} & 15 & 0.5956 & 0.0906 & 0.0073 & 0.7478 & 0.7801 & 0.7970 & 0.7954 & 0.7986 & 23017 \\
\texttt{pysteps} & 30 & 0.6489 & 0.1001 & 0.0096 & 0.7265 & 0.7404 & 0.7596 & 0.7903 & 0.7313 & 21324 \\
\texttt{pysteps} & 45 & 0.7006 & 0.1107 & -0.0111 & 0.6681 & 0.5968 & 0.6216 & 0.7865 & 0.5138 & 15133 \\
\texttt{pysteps} & 60 & 0.7837 & 0.1226 & -0.0618 & 0.4618 & 0.3821 & 0.4084 & 0.7491 & 0.2808 & 8727 \\
\end{longtable}
}

\begin{thebibliography}{99}

\bibitem{messer2006}
H.~Messer, A.~Zinevich, and P.~Alpert.
\newblock Environmental monitoring by wireless communication networks.
\newblock \emph{Science}, 312(5774):713, 2006.

\bibitem{goldshtein2009}
O.~Goldshtein, H.~Messer, and A.~Zinevich.
\newblock Rain rate estimation using measurements from commercial telecommunications links.
\newblock \emph{IEEE Transactions on Signal Processing}, 57(4):1616--1625, 2009.

\bibitem{pulkkinen2019pysteps}
S.~Pulkkinen, D.~Nerini, A.~A. P\'erez Hortal, C.~Velasco-Forero, A.~Seed, U.~Germann, and L.~Foresti.
\newblock pySTEPS: An open-source Python library for probabilistic precipitation nowcasting.
\newblock \emph{Geoscientific Model Development}, 12(10):4185--4219, 2019.

\bibitem{vaswani2017attention}
A.~Vaswani, N.~Shazeer, N.~Parmar, J.~Uszkoreit, L.~Jones, A.~N. Gomez, {\L}.~Kaiser, and I.~Polosukhin.
\newblock Attention is all you need.
\newblock In \emph{Advances in Neural Information Processing Systems (NeurIPS)}, 2017.

\bibitem{brunton2016sindy}
S.~L. Brunton, J.~L. Proctor, and J.~N. Kutz.
\newblock Discovering governing equations from data by sparse identification of nonlinear dynamical systems.
\newblock \emph{Proceedings of the National Academy of Sciences}, 113(15):3932--3937, 2016.

\bibitem{gu2023mamba}
A.~Gu and T.~Dao.
\newblock Mamba: Linear-time sequence modeling with selective state spaces.
\newblock \emph{arXiv preprint arXiv:2312.00752}, 2023.

\bibitem{ituR838}
International Telecommunication Union.
\newblock \emph{Recommendation ITU-R~P.838-3: Specific attenuation model for rain for use in prediction methods}.
\newblock ITU, Geneva, Switzerland, 2005.

\bibitem{overeem2013}
A.~Overeem, H.~Leijnse, and R.~Uijlenhoet.
\newblock Country-wide rainfall maps from cellular communication networks.
\newblock \emph{Proceedings of the National Academy of Sciences}, 110(8):2741--2745, 2013.

\bibitem{marshallpalmer1948}
J.~S. Marshall and W.~M. Palmer.
\newblock The distribution of raindrops with size.
\newblock \emph{Journal of Meteorology}, 5(4):165--166, 1948.

\bibitem{eshel2021spatial}
A.~Eshel, H.~Messer, H.~Kunstmann, P.~Alpert, and C.~Chwala.
\newblock Quantitative analysis of the performance of spatial interpolation methods for rainfall estimation using commercial microwave links.
\newblock \emph{Journal of Hydrometeorology}, 22(4):831--843, 2021.
\newblock DOI: \texttt{10.1175/JHM-D-20-0164.1}.

\bibitem{andersson2023openmrg}
J.~C.~M. Andersson, J.~Olsson, R.~van~de~Beek, and J.~Hansryd.
\newblock openMRG: Open data from microwave links, radar, and gauges for rainfall quantification in Gothenburg, Sweden.
\newblock \emph{Earth System Science Data}, 2023.

\bibitem{chwala2019cml}
C.~Chwala and H.~Kunstmann.
\newblock Commercial microwave link networks for rainfall observation: Assessment of the current status and future challenges.
\newblock \emph{WIREs Water}, 6(2):e1337, 2019.

\bibitem{habi2021pynncml}
H.~V. Habi and H.~Messer.
\newblock Wet-dry classification using LSTM and commercial microwave links.
\newblock In \emph{Proceedings of the IEEE Sensor Array and Multichannel Signal Processing Workshop (SAM)}, 2018.

\bibitem{devos2019}
L.-W. de~Vos, A.~Overeem, H.~Leijnse, and R.~Uijlenhoet.
\newblock Rainfall estimation accuracy of a nationwide instantaneously sampling commercial microwave link network: Error dependency on known characteristics.
\newblock \emph{Journal of Hydrometeorology}, 20(6):1101--1115, 2019.

\end{thebibliography}

\end{document}
